\clearpage
\part*{Part V. The Matrix, The Standard Model At The Six Genres}
\label{atlaspart:05-gauge-structure}
\addcontentsline{toc}{part}{Part V. The Matrix, The Standard Model At The Six Genres}
\section*{Introduction to Part V}\label{narr:part:IV}
\addcontentsline{toc}{section}{Introduction to Part V}
The Standard Model of the literature is a gauge theory with a group, a set of representations, a charge assignment, conservation laws and anomaly conditions. Part V asks which of these structures the construction produces from the known clause of the axiom and the six genres of the link, and in what order. The answer is a thread that begins with gauge invariance read as the known clause, passes through the group, the representations, the weights and the charge, and ends with the conservation laws, the Ward identity and the table of the seventy-five entries of the model with their standing. Along that thread the Part is explicit about what is a theorem, what is conditional and what is open, because this is the Part where the names of the literature press hardest on the construction.

The order of the ten chapters is the order of a textbook on gauge theory read backward from its premises. A textbook begins with a group and a set of fields and derives invariance; this Part begins with invariance, which is the known clause of the axiom read at one link, and derives the group from it. The representations are then read as a census of open ends, the weights and the charge as functionals at those ends, and the sign as the side of an acting site, so that each chapter reads the next structure off the one ahead of it. The conservation chapters that follow are where the construction is most unlike the textbook: conservation is a kernel of a move matrix and not a symmetry of an action, and the Noether chapter records that the two notions part at every sink. The Ward identity and the word table are the last readings at one link, and the table of the seventy-five entries closes the Part by printing the standing of every name. A reader who wants the two theorems that carry the Part should read chapters 4-1 and 4-2 first and treat everything from chapter 4-3 onward as a sequence of readings whose standing is printed inside each box.

Chapter 4-1 starts from the definition of known. Theorem~64.1 reads gauge invariance as the known clause, and Theorem~64.2 reads the gauge group as the face group of the tether. The known clause is the matching of basis directions across pairs, and every readout of the one qutrit-qubit link is invariant under every matching of the identification group as a corollary of that clause, the group re-met as exact objects at the fourth-root monomial groups of the orders four, thirty-two, three hundred eighty-four and six thousand one hundred forty-four and at the whole group at the Lie dimensions four and nine, Theorem~64.1. The face group of the tether, the image of the local unitaries of the two ends on the link, is the connected determinant slice with a Lie algebra of dimension twelve, a derived algebra of dimension eleven and a one-dimensional center, isomorphic to the quotient of the product of the three factors by the cyclic group of order six through an explicit covering, and the uniqueness cell at the dimension twelve and the rank four picks the qutrit-qubit type, Theorem~64.2. Chapter 4-2 then obtains the group as a corollary of one construction in Theorem~65.1, with the hub pair as one qudit of dimension four in Theorem~65.2. The one construction is the stabilizer of the Hermitian unit inside the isometries of the split, the same object chapter 3-2 used for the rotations, applied at the three end kinds of the listed webs: at the hub-leg qubit copies it gives dimension three and rank one with the lift center of order two, the weak factor, and at the qutrit ends open at the hub leg dimension eight and rank two with the lift center of order three, Theorem~65.1. At the ququart hub ends the hub pair read as one qudit of dimension four yields the Lie algebra of dimension fifteen and rank three with the center cyclic of order four, the adopted carrier phase datum a generator of that center, and the hub-pair weights at the two qubit copies with the u(1) weight at the dock slice equal the doublet of the Higgs entry by entry, so that the Higgs representation is an output of the hub-pair weights, Theorem~65.2. These two chapters are the theorems of the Part; what follows them is a sequence of open and conditional results that the reader should take at their stated standing.

Chapter 4-3 turns to the representations. Open result~66.1 states them from the open-end census, Open result~66.2 reads chirality as the dock count of the read, Theorem~66.4 reads the carriers as the closed reads of their type faces, and Open result~66.5 reads the roster as the sinks of the move graph; two further counts of the representation census, Open result~67.1 and Open result~68.1, stand beside them. The representation content is the open-end census of a body, an output at every web and every sink, and the sixteen rows of the model are reached at the wide union of the readings and at no one reading, which is the parting, Open result~66.1. The electron reaches both of its rows, one dock entered the left doublet and both docks the right singlet, so the dock-count reading of chirality is derived there while the isospin sign from the orientation of the open end parts at the neutral body, Open result~66.2. Every carrier of the list is the closed read of its type face, the W at the (2,2) type, the gluon at the (3,3) type and the photon the open (2,3) pair, Theorem~66.4, and within the charge fibers the sinks of the move graph equal the roster at thirty-five of the sixty-four listed webs, the first differing web a three-quark body, Open result~66.5. The two censuses beside them count the saturated sinks, sixteen thousand nine hundred sixty-two at the second pass, where the hub count outputs two generations at the lepton-form sinks, Open result~67.1, and twenty-two thousand three hundred two at the union with the tau-entry sinks, where the generation count is three at both readings and the parting of the second pass dissolves, Open result~68.1. Chapter 4-4 reads the Cartan weights and the charge, Open result~72.1, with the charge as a centralizer weight at the dock face, Open result~72.3, and the signed open-end count under the persistence matching, Theorem~72.4. The weights are assigned at every open end of every listed web and every sink, and they meet the dart at four of the twenty-two listed webs carrying an open end, Open result~72.1. The charge lies inside the centralizer of the dock-face stabilizer and outside the centralizer of the construction whole, and the dart re-read in the signed-dart-end basis equals that weight at seventeen of the twenty-two listed webs at the dock-rule face and at fifteen at the hub-leg face, Open result~72.3. Under the persistence matching the signed open-end count at the unit weight is conserved at all but eight hundred thirty-five of the move instances of the census, every changing instance a re-matching of two ends of one type, so that the conservation is a theorem of the joins and an open result of the re-matchings, Theorem~72.4. Chapter 4-5 supplies the sign: Theorem~73.1 obtains the sign from the side at the hub reached through its qutrit, at all five hubs of the ladder so entered and at every hub of the three lepton-seeded sinks, with the slot rule shown not to be the sink orientation read at the cut end.

Chapters 4-6 to 4-8 are the conservation chapters. Theorem~74.1 gives the conserved functionals of the full move set, Theorem~74.4 the junction count at the moves, and Open result~75.1, as an open result, the kernel as the center of the Weyl algebra at the sinks. The conserved functionals of the full move set are the left kernel of the move matrix, trivial over the grade vector at both faces, with the lepton number as the dark-triangle count broken at a freed-at-cut column and the dart charge at the pure open-leg column, Theorem~74.1. The register at which energy and momentum conserve across every junction is the dock, a corollary of the dock closure, and the momentum transfer composes to the zero signed step along the direction register at every record, Theorem~74.2. Baryon number read as the unsigned junction count is conserved at every interaction and moved by the non-interactions alone, Theorem~74.4. At the seven derived capped sinks the Weyl algebra of the cap records has a center of dimension one at every sink, and the commutant of the admitted group moves parts from that center at every one of the seven, first at the proton sink with a commutant of dimension twelve against the center of one, Open result~75.1. Chapter 4-7 reads the Noether theorem at the seeded sinks, Open result~76.1, and the conserved charges as the center census, Open result~76.2, both open. Noether at the sinks is the kernel of the move matrix and not the commutant of the group moves nor the Cartan centralizer, which act on no count, the three compared cell by cell at the fifty-three seeded sinks agreeing at no sink, Open result~76.1. Chapter 4-8 reads the Ward identity from conservation as one link, Theorem~78.1, the anomaly at the body part, Theorem~78.3, and the anomaly with every body at its own sink, Theorem~79.1. The identity holds at every one of the forty-one readings of the records the definitions admit, so the adopted reading is derived, Theorem~78.2, and the body-part anomaly is defined with the orders separated, at order one the unit itself with the sign plus, at order two an exact element of the field with the sign plus and a nonzero cross term, and the order-three term zero, Theorem~78.3.

Chapter 4-9 states the one rule of a read and the word table, Open result~84.1, with the passage web of the cross read against the W word, Open result~84.2, and the Z read as the two-line read of the W word, Open result~84.3, every one an open result that Part VIII takes up again. Applied at the sixty seeded sinks with every end lettered, the rule reaches nine of the sixteen words at the one face under the adopted convention, thirteen at the union of three conventions and the sixteen at the wide union alone, Open result~84.1. Chapter 4-10 closes with the table of the seventy-five entries (in chapter~\ref{mod:IV-10-sm-table-the-75-entries-with-their-status}, among the instruments), which is the place to look up the standing of any entry of the model as this compendium reads it, every grade read from a derived entry, the four hundred seventy-two entries of the corpus at the grades derived, named, silent, unmapped, gloss and open with the dependency graph acyclic.

In the language of the literature, this Part names gauge invariance, the gauge group, the representations, chirality, the Cartan weights, the electric charge, the Noether theorem, the Ward identity, the anomaly and the Standard Model itself. The correspondence is by name at each. The boxes establish finite structure: a face group, a census of open ends, weights read at a dock face, functionals in the kernel of a move matrix.

The dependency lines of this Part reach into Parts I and III and forward into Part VIII, since the doublets of chapter 6-3 enter the weights and the sign. Gauge invariance rests on the axiom chapter, the link and the records, the group on the rotation construction of Part IV, the representations on the moves and the saturation, and the anomaly on the measurement web of Part VII. A reader following the derivation graph will find Part V the most tightly knotted region of the compendium, and the path table at the front gives the shortest chain from the axiom to its first theorem. The reader should also notice that the two words the textbook treats as one, invariance and conservation, are here two different objects, a clause of the axiom and a kernel of a matrix, whose relation is itself an open result of chapter 4-7.
\section{Gauge invariance as the known clause}\label{mod:IV-1-gauge-invariance-as-the-known-clause}
\chdesig{5-1}{4-1}{IV-1-gauge-invariance-as-the-known-clause}{IV-1-gauge-invariance-as-the-known-clause}
\label{narr:open:4-1}Gauge invariance enters as a reading of the definition of known. Theorem~64.1 reads gauge invariance as the known clause and Theorem~64.2 reads the gauge group as the face group of the tether, both resting on the axiom chapter, the link and the records.
\begin{chtheorem}[Gauge invariance as the known clause]\label{box:IV-1-gauge-invariance-as-the-known-clause:1}
This chapter establishes that gauge invariance is the known clause (the known clause: the matching of basis directions across pairs) read at the one link. Every readout of the one qutrit-qubit link is invariant under every matching of the identification group as a corollary of that clause: the group is re-met as exact objects, the fourth-root monomial groups at the orders four, thirty-two, three hundred eighty-four and six thousand one hundred forty-four, the whole group at the Lie dimensions four and nine, and the identification algebra at the six link types with the one readable type carrying the acting centre; the one hundred seventy-two records are rebuilt and every readout is invariant on the state under every matching, every Lie generator, the Cayley unitaries, the chain unitaries and the Fourier elements, the corollary standing with no posed step. The group is named as the gauge group of the Standard Model by correspondence.
\end{chtheorem}
\chstanding{theorem}
\chdeps{Definitions 1.1 to 1.6, Theorems 2.1 to 2.3, Open results 2.4 and 2.5, Theorem 2.6, Open results 2.7 and 2.8 and Theorems 2.9 and 2.10, Theorems 9.1 to 9.7 and Open result 9.8 and Theorems 18.1 to 18.4 and Open results 18.5 and 18.6}{Theorems 65.1 and 65.2, Open results 66.1 to 66.3, Theorem 66.4 and Open result 66.5, Open results 72.1 to 72.3, Theorem 72.4 and Open result 72.5, Theorems 74.1 to 74.4 and Open result 74.5, Open results 76.1 to 76.3, Open results 84.1 to 84.3 and Theorem 163.1, Open result 163.2 and Theorem 163.3}
\begin{chtheorem}[The gauge group as the face group of the tether]\label{box:IV-1-gauge-invariance-as-the-known-clause:2}
The group itself: the face group of the qutrit-qubit tether, the image of the local unitaries $U(2)$ x $U(3)$ on the link, is the connected determinant slice $S(U(2)$ x $U(3))$ with Lie algebra of dimension twelve, derived algebra of dimension eleven and a one-dimensional centre spanned by the generator of the identification-algebra result, isomorphic to $(SU(2)$ x $SU(3)$ x $U(1))/\mathbb{Z}_6$ by an explicit covering whose kernel is the cyclic group of order six; the uniqueness cell at the dimension twelve and the rank four picks the qutrit-qubit type at both, and the sixteen-row law of the dart holds at every row.
\end{chtheorem}
\chstanding{theorem}
\chdeps{Definitions 1.1 to 1.6, Theorems 2.1 to 2.3, Open results 2.4 and 2.5, Theorem 2.6, Open results 2.7 and 2.8 and Theorems 2.9 and 2.10, Theorems 9.1 to 9.7 and Open result 9.8 and Theorems 18.1 to 18.4 and Open results 18.5 and 18.6}{Theorems 65.1 and 65.2, Open results 66.1 to 66.3, Theorem 66.4 and Open result 66.5, Open results 72.1 to 72.3, Theorem 72.4 and Open result 72.5, Theorems 74.1 to 74.4 and Open result 74.5, Open results 76.1 to 76.3, Open results 84.1 to 84.3 and Theorem 163.1, Open result 163.2 and Theorem 163.3}
The known clause names what every read of a qudit fixes jointly whichever way the basis directions are matched across its pairs; the identification group is the set of those matchings, built as exact objects at the one link, and the readout of a record is a trace against the twirl over that group, which is fixed by every element it averages, so the sum is re-indexed and every readout is invariant; the records and the readouts are recomputed by two construction-disjoint engines and agree at every cell. The local unitaries of the two ends act on the link, the kernel of that action on the determinant slice is computed as a group, the image is the slice with its Lie algebra decomposed into its ideals, and the covering is written explicitly with its kernel enumerated; the sixteen-row law is checked by integer arithmetic.

\begin{chbridge}
A gauge theory is a field theory whose Lagrangian is invariant under a group of transformations acting independently at every point of spacetime; the gauge group of the Standard Model is $SU(3)_c$ x $SU(2)_L$ x $U(1)_Y$, the group acting faithfully on the matter representations being its quotient by a cyclic subgroup of order six of the centre \cite{YangMills1954,Weinberg1967,Baez2010}. The derivation supplies the identification group of the one link with its invariant readouts and the face group of the tether as exact objects here, and gauge invariance and the gauge group are named by correspondence, the representations standing at chapter 4-3 and the charge at chapter 4-4.
\end{chbridge}
\begin{chproof}[finite and census]
The premises are the axiom with the definition \emph{known} (chapter 1-1, the text), the admissible matchings as the whole unitary group at each end (chapter 3-6, a theorem), the identification algebra at the (2,3) tether (a theorem placed at no chapter of this compendium), the explicit two-point function with its records (chapter 5-1, a theorem) and the Born count as the readout of a record (chapter 2-2, a theorem). The chain back to the axiom runs through chapter 1-1 and the framework chapter. The procedure. The input is the one (2,3) link with its one hundred seventy-two records and the identification group of its two ends, built as exact objects, the fourth-root monomial groups at the orders four, thirty-two, three hundred eighty-four and six thousand one hundred forty-four, the whole group at the Lie dimensions four and nine with its thirteen generators, and the identification algebra at the six link types with the one readable type carrying the acting center. For every record the readout is the trace of its projector against the twirled state over the register dimension,
\begin{equation}
\langle P\rangle = \frac{\mathrm{tr}\,\big(M\,T(P)\big)}{\dim R},\qquad T(P) = \frac{1}{|G|}\sum_{g\in G} g\,P\,g^{-1} ,
\end{equation}
and the test applied at every record and every element is whether the readout is unmoved by the element, the elements being the three hundred eighty-four qutrit-end matchings, the thirty-two qubit-end matchings, the thirteen Lie generators (the derivative of the trace along each one), the thirteen Cayley unitaries, the two chain unitaries and the two Fourier elements. The aggregation counts the records unmoved under each family. The totals at the checkpoints at the reading taken, the whole-group twirl, unmoved at one hundred seventy-two of one hundred seventy-two records under every family, the first moved record none. So every readout of the one link is invariant under every matching as a corollary of the known clause with no posed step, which is the theorem, the group named as the gauge group of the Standard Model by correspondence. The smallest instance on the page is one record under one matching, the twirl is fixed by the element it averages, $g\,T(P)\,g^{-1} = T(P)$, so the trace is re-indexed and the readout unmoved. At the faces beside, the finite-group twirl is moved by the Lie generators at two records and by the Fourier elements at fourteen, and the bare untwirled state is moved by the three hundred eighty-four matchings at nine records, printed beside.

The counter-tests of the derivation stood at their pinned letters and every control held. The records and the readouts were recomputed by two engines built apart with no shared line and agree at every cell.
\chcert{384}
\end{chproof}
\par\medskip\noindent{\itshape With invariance read, the next chapter obtains the group itself from one construction.}\label{narr:close:4-1}\par
\section{The gauge group as a corollary of one construction}\label{mod:IV-2-gauge-group-as-a-corollary-of-one-construction}
\chdesig{5-2}{4-2}{IV-2-gauge-group-as-a-corollary-of-one-construction}{IV-2-gauge-group-as-a-corollary-of-one-construction}
\label{narr:open:4-2}Theorem~65.1 obtains the gauge group as a corollary of one construction, and Theorem~65.2 reads the hub pair as one qudit of dimension four, the two theorems on which the rest of the Part rests.
\begin{chtheorem}[The gauge group as a corollary of one construction]\label{box:IV-2-gauge-group-as-a-corollary-of-one-construction:1}
This chapter establishes that the gauge group is a corollary of one construction. The stabiliser of the Hermitian unit inside the isometries of the split, spanned by the adjoint matrices of the derived directions, applied at the three end kinds of the listed webs, yields the three factors as exact objects: at the hub-leg qubit copies the register of dimension two gives dimension three and rank one with the lift centre of order two, the weak factor; at the qutrit ends open at the hub leg the register of dimension three gives dimension eight and rank two with the lift centre of order three, the colour factor; and at the dock qubit ends the $u(1)$ of the unit has the slice representative equal to the centre generator of the group result, while the stabiliser at the dock register is the rotation group, spatial and compared with no factor; the kernel of the covering is re-met from the three register centres. The reading beside, the stabiliser inside the isometries whole, parts from its factor first at the hub-leg copy.
\end{chtheorem}
\chstanding{theorem}
\chdeps{Definitions 1.1 to 1.6, Theorems 9.1 to 9.7 and Open result 9.8, Theorems 64.1 and 64.2 and Open results 241.1 and 241.2, Theorems 241.3 and 241.4 and Open results 241.5 to 241.7}{Open results 66.1 to 66.3, Theorem 66.4 and Open result 66.5, Open results 72.1 to 72.3, Theorem 72.4 and Open result 72.5, Theorem 73.1 and Open results 73.2 and 73.3, Open result 75.1, Open results 76.1 to 76.3, Theorem 163.1, Open result 163.2 and Theorem 163.3, Theorems 200.1 to 200.3 and Open result 200.4 and Open result 236.1}
\begin{chtheorem}[The hub pair as one qudit of dimension four]\label{box:IV-2-gauge-group-as-a-corollary-of-one-construction:2}
At the ququart hub ends, the hub pair read as one qudit of dimension four by the Axiom, the same construction yields the Lie algebra of dimension fifteen and rank three with the lift $SU(4)$ and the centre cyclic of order four; the adopted carrier phase datum, the fourth root, is a generator of that centre, so the datum is derived at the hub centre; and the hub-pair weights at the two qubit copies with the $u(1)$ weight at the dock slice equal the doublet of the Higgs entry by entry, so the Higgs representation is an output of the hub-pair weights, the readings beside parting at the centre or at a weight.
\end{chtheorem}
\chstanding{theorem}
\chdeps{Definitions 1.1 to 1.6, Theorems 9.1 to 9.7 and Open result 9.8, Theorems 64.1 and 64.2 and Open results 241.1 and 241.2, Theorems 241.3 and 241.4 and Open results 241.5 to 241.7}{Open results 66.1 to 66.3, Theorem 66.4 and Open result 66.5, Open results 72.1 to 72.3, Theorem 72.4 and Open result 72.5, Theorem 73.1 and Open results 73.2 and 73.3, Open result 75.1, Open results 76.1 to 76.3, Theorem 163.1, Open result 163.2 and Theorem 163.3, Theorems 200.1 to 200.3 and Open result 200.4 and Open result 236.1}
The end kinds of the listed webs are located at their copies, the derived directions of each register are listed, the stabilizer of the Hermitian unit inside the isometries of the split is spanned by their adjoint matrices and its dimension, rank and lift center computed as exact objects, and the three results are compared factor by factor with the gauge group result and with the rotation group result; the hub pair is read as one qudit of dimension four, the construction is run at the one-pair hubs and at the ends of the listed ququart-ququart links, the center of the resulting algebra is enumerated at the scalars, the carrier phase datum is compared with its generators, and the hub-pair weights are read at the two qubit copies and compared with the Higgs doublet entry by entry.

\begin{chbridge}
The gauge group of the Standard Model is a product of three factors, the colour factor $SU(3)$ of dimension eight and rank two with centre $\mathbb{Z}_3$, the weak factor $SU(2)$ of dimension three and rank one with centre $\mathbb{Z}_2$, and the hypercharge factor $U(1)$ \cite{Weinberg1967,GeorgiGlashow1974}; a compact Lie group is determined up to covering by its Lie algebra, and the stabiliser of a vector under a unitary group is a closed subgroup whose Lie algebra is the annihilator of that vector \cite{Hall2015,Humphreys1972}; $su(4)$ has dimension fifteen and rank three with centre $\mathbb{Z}_4$, and the Higgs field is a doublet of $SU(2)$ of hypercharge one half. The derivation supplies the stabilizer of the unit at the three end kinds as exact Lie algebras with their lift centers here, and the gauge group, the rotation group, and the Higgs doublet are named by correspondence, the representations standing at chapter 4-3 by the open-end census alone.
\end{chbridge}
\begin{chproof}[finite and exact]
The premises are the link with its six types and the listed webs with their end kinds (chapter 1-2, a theorem), the stabilizer construction at the one link (chapter 3-2, a theorem) and the gauge group derived at the identification algebra (chapter 4-1, a theorem). The chain back to the axiom runs through chapter 4-1 and the framework chapter. The procedure. The input is the three end kinds of the listed webs located at their copies, six hundred fifty-seven hub-leg qubit copies, two hundred nineteen qutrit ends open at the hub leg and twenty-five dock qubit ends, with the derived directions of each register. At each register of dimension $N$ the one construction is applied, the stabilizer of the Hermitian unit inside the isometries of the split, spanned by the adjoint matrices of the derived directions, and its dimension, rank and lift center are computed as exact objects. The test applied at each kind is whether the algebra obtained equals the factor of the gauge group at that kind, and the aggregation counts the kinds equal. The exact part. From a register of dimension $N$, by the span of the adjoint matrices of its derived directions, the traceless Hermitian matrices, we obtain
\begin{equation}
\dim = N^{2}-1,\qquad \operatorname{rank} = N-1,\qquad \text{the lift center } \mathbb{Z}/N ,
\end{equation}
so at the hub-leg qubit copies the register of dimension two gives dimension three, rank one and the center of order two, the weak factor. At the qutrit ends the register of dimension three gives dimension eight, rank two and the center of order three, the color factor, and at the dock qubit ends the $u(1)$ of the unit has the slice representative
\begin{equation}
(3,\,3,\,-2,\,-2,\,-2) ,
\end{equation}
equal to the center generator of the gauge-group result, with the stabilizer at the dock register the rotation group $so(3)$, spatial and compared with no factor. The kernel of the covering is re-met from the three register centers at the six triples of the gauge-group result with the generator $(3,2,1)$, the cyclic group of order six. The totals at the checkpoints, the three kinds equal to their factors at three of three, so the gauge group is a corollary of the one construction at three end kinds, which is the theorem. At the reading beside, the stabilizer inside the isometries whole, the hub-leg copy gives the connected rotation group with the scalar center of order one against the weak factor with its center of order two, and the qutrit end gives an algebra of dimension twenty-eight and rank four against dimension eight and rank two, the first parting at the hub-leg copy. The smallest instance on the page is a register of dimension two, its three derived directions are the traceless Hermitian matrices of a qubit, their adjoint matrices span an algebra of dimension three and rank one, and its lift carries the center of order two.

The counter-tests struck the Hermitian unit from the split at every register, which left the stabilizer undefined and moved the cells of the reading beside. read the qutrit register at dimension two, which parted the color comparison and moved the deciding cell, and struck the text at its gate, which left the gate unmet. Every control read the same at both engines but one, which counted three rows the parent itself excludes, disclosed. The algebras were reproduced at a second engine, equal at every kind.
\chcert{901}
\end{chproof}
\par\medskip\noindent{\itshape The group in hand, the next chapter counts its representations at the open ends.}\label{narr:close:4-2}\par
\section{The representations from the open-end census}\label{mod:IV-3-representations-from-the-open-end-census}
\chdesig{5-3}{4-3}{IV-3-representations-from-the-open-end-census}{IV-3-D543954-the-wave-as-the-carrier-advance-and-the-removing-moves}
\label{narr:open:4-3}The representations are read from a census of open ends. Open result~66.1 states them, Open result~66.2 reads chirality as the dock count of the read, Open result~66.3 gives the census at the sinks and the words reached, Theorem~66.4 reads the carriers as the closed reads of their type faces, and Open result~66.5 reads the roster as the sinks of the move graph; the second and third counts of the census, Open result~67.1 and Open result~68.1, stand beside them as open results.
\begin{chwall}[The representations from the open-end census]\label{box:IV-3-representations-from-the-open-end-census:1}
This chapter reads the representation content as the open-end census of a body, an output at every web and every sink, and finds the sixteen rows reached at the wide union of the readings and at no one reading. At the reading of the listed end lines the eight states of the four bodies read equal to the sixteen rows at two, and at the transport reading of the open end at two as well, the two readings reaching complementary pairs of rows; at the lepton ladder the muon and tau rows read equal at two of four and, with the definition of open in place of the adopted reading, at none, parting first at the muon.
\end{chwall}
\chstanding{open result}
\chdeps{Definitions 1.1 to 1.6, Theorems 9.1 to 9.7 and Open result 9.8, Theorems 18.1 to 18.4 and Open results 18.5 and 18.6, Theorems 28.1 and 28.2, Open results 28.3 and 28.4 and Theorem 28.5, Theorems 64.1 and 64.2, Theorems 65.1 and 65.2, Open results 241.1 and 241.2, Theorems 241.3 and 241.4 and Open results 241.5 to 241.7 and Open results 245.1 and 245.2, Theorem 245.3, Open results 245.4 to 245.8 and Theorem 245.9}{the theorems of Chapters 5-3 to 5-4, 5-6, 5-8 to 5-9, 5-11, 8-3, 11-1, 14-6}
\begin{chwall}[Chirality as the dock count of the read]\label{box:IV-3-representations-from-the-open-end-census:2}
Read as the set of reads, the electron reaches both of its rows, one dock entered the left doublet and both docks the right singlet, so the dock-count reading of chirality is derived there, while the isospin sign from the orientation of the open end parts at the neutral body; the closed reads of the W and the gluon reach their adjoints and the photon the singlet, but the Higgs attachment read opens no qubit end against the Higgs doublet.
\end{chwall}
\chstanding{open result}
\chdeps{Definitions 1.1 to 1.6, Theorems 9.1 to 9.7 and Open result 9.8, Theorems 18.1 to 18.4 and Open results 18.5 and 18.6, Theorems 28.1 and 28.2, Open results 28.3 and 28.4 and Theorem 28.5, Theorems 64.1 and 64.2, Theorems 65.1 and 65.2, Open results 241.1 and 241.2, Theorems 241.3 and 241.4 and Open results 241.5 to 241.7 and Open results 245.1 and 245.2, Theorem 245.3, Open results 245.4 to 245.8 and Theorem 245.9}{the theorems of Chapters 5-3 to 5-4, 5-6, 5-8 to 5-9, 5-11, 8-3, 11-1, 14-6}
\begin{chwall}[The open-end census at the sinks and the words reached]\label{box:IV-3-representations-from-the-open-end-census:3}
At the sinks the open-end census is an output at every one, yet the listed muon lies among the sinks of neither closure and the anomaly sums over the words reached are nonzero at the adopted face; at the hub-leg reading six of the sixteen words are reached; every word is reached at some seeded sink at the wide union over every carried face and both conventions of the weak component, while no one face under one convention reaches the sixteen; the reach is not complete at the one reading of the seeded sinks; and the colour rows are outputs at every quark-bearing sink but one, the two-dock down body.
\end{chwall}
\chstanding{open result}
\chdeps{Definitions 1.1 to 1.6, Theorems 9.1 to 9.7 and Open result 9.8, Theorems 18.1 to 18.4 and Open results 18.5 and 18.6, Theorems 28.1 and 28.2, Open results 28.3 and 28.4 and Theorem 28.5, Theorems 64.1 and 64.2, Theorems 65.1 and 65.2, Open results 241.1 and 241.2, Theorems 241.3 and 241.4 and Open results 241.5 to 241.7 and Open results 245.1 and 245.2, Theorem 245.3, Open results 245.4 to 245.8 and Theorem 245.9}{the theorems of Chapters 5-3 to 5-4, 5-6, 5-8 to 5-9, 5-11, 8-3, 11-1, 14-6}
\begin{chconditional}[The carriers as the closed reads of their type faces]\label{box:IV-3-representations-from-the-open-end-census:4}
Conditional on the join face as posed, and on the results these theorems rest on at their own standing.
The vertex census of the qutrit-qubit link is a function of its two end dimensions alone; the adopted W is built and equals the cross of the listed electron web, while no closed Z web is admitted at the list rules; the spin of every carrier is read by name at the read link, conditional on the join face as posed; every carrier of the list is the closed read of its type face, the W at (2,2), the gluon at (3,3), the photon the open (2,3) pair; and the derived carrier reads differ from the listed words at the signature, the degeneracy of the listed words dissolved.
\end{chconditional}
\chstanding{conditional theorem}
\chdeps{Definitions 1.1 to 1.6, Theorems 9.1 to 9.7 and Open result 9.8, Theorems 18.1 to 18.4 and Open results 18.5 and 18.6, Theorems 28.1 and 28.2, Open results 28.3 and 28.4 and Theorem 28.5, Theorems 64.1 and 64.2, Theorems 65.1 and 65.2, Open results 241.1 and 241.2, Theorems 241.3 and 241.4 and Open results 241.5 to 241.7 and Open results 245.1 and 245.2, Theorem 245.3, Open results 245.4 to 245.8 and Theorem 245.9}{the theorems of Chapters 5-3 to 5-4, 5-6, 5-8 to 5-9, 5-11, 8-3, 11-1, 14-6}
\begin{chwall}[The roster as the sinks of the move graph]\label{box:IV-3-representations-from-the-open-end-census:5}
The roster is laid as the sinks of the move graph and of the one-act moves; within the charge fibres the sinks equal the roster at thirty-five of the sixty-four listed webs, the first differing web a three-quark body; the readable sinks within the fibres are the thirteen zero-sum webs and the gluon word; the roster is not the set of join-irreducible webs; and the office-bearing rule is necessary at the thirteen bodies and selects nothing among the listed candidates.
\end{chwall}
\chstanding{open result}
\chdeps{Definitions 1.1 to 1.6, Theorems 9.1 to 9.7 and Open result 9.8, Theorems 18.1 to 18.4 and Open results 18.5 and 18.6, Theorems 28.1 and 28.2, Open results 28.3 and 28.4 and Theorem 28.5, Theorems 64.1 and 64.2, Theorems 65.1 and 65.2, Open results 241.1 and 241.2, Theorems 241.3 and 241.4 and Open results 241.5 to 241.7 and Open results 245.1 and 245.2, Theorem 245.3, Open results 245.4 to 245.8 and Theorem 245.9}{the theorems of Chapters 5-3 to 5-4, 5-6, 5-8 to 5-9, 5-11, 8-3, 11-1, 14-6}
Every end of every web and sink is lettered under the reading at hand, the state of the body is expanded and compared row by row with the sixteen, and the first parting row is named at every reading; the carriers are read as the closed reads of their type faces with the records surviving both signed exchanges counted as unordered pairs of far-end letters; the orbit count of the qutrit-qubit link is written as a polynomial in the two end dimensions; and the sinks within the charge fibers are compared with the roster web by web.

\begin{chbridge}
The matter content of the Standard Model at one generation is a list of representations of $SU(3)_c$ x $SU(2)_L$ x $U(1)_Y$: the lepton doublet $(1,2)_{-1/2}$ and singlet $(1,1)_{-1}$, the quark doublet $(3,2)_{1/6}$ and singlets $(3,1)_{2/3}$ and $(3,1)_{-1/3}$, with $Q = T_3 + Y$ \cite{Weinberg1967,FritzschGellMannLeutwyler1973}; the carriers transform in the adjoint representations and the Higgs in the doublet $(1,2)_{1/2}$ \cite{Weinberg1967}; the gauge anomalies cancel at every generation \cite{Bouchiat1972,GeorgiGlashow1974}. The derivation supplies the open-end census of a body with its weights and the closed reads of the type faces here, and the representations, chirality, the carriers, and the Higgs doublet are named by correspondence, the sixteen rows compared with being an input to no construction.
\end{chbridge}
\begin{chproof}[finite and census]
The premises are the listed webs with their end lines and the four bodies (chapter 1-5, a theorem), the Cartan weights at the open ends (chapter 4-4, an open result read at its standing), the gauge group with its sixteen rows as a comparison target and no input (chapter 4-1, a theorem) and the dart charge at the dock face (an open result placed at no chapter of this compendium, read at its standing). The chain back to the axiom runs through chapter 1-5 and the framework chapter. The procedure. The input is the eight states of the four bodies at the reading taken, the end lines of the list as the open ends, the ket copy entered by the read line and washed to the pair ket on the state, the bra copy its conjugate, the dock copies closed at the dock face and the first-order face. Every open end is lettered under the reading, the state of each body is expanded into rows with its weak and hypercharge weights in sixths and its color multiplicity, and the rows are compared one by one with the sixteen rows at the two assignments of the ket read to the left or the right rows. The test applied is equality of a derived row with a listed row, and the aggregation counts the rows equal and names the first parting row, repeated at the sixty readings of the grid at both assignments. The totals at the checkpoints, the eight states expand to fourteen rows against the sixteen. At the assignment of the ket read to the left rows the rows equal at two of sixteen, the two electron rows, the doublet at the odd-leg qubit with the weight
\begin{equation} \adjustbox{max width=\dimexpr\linewidth-3em\relax}{$\displaystyle \begin{gathered} (6T_{3},\ 6Y) = (-3,\,-3) \ \text{at the left row}, \\ (0,\,-6) \ \text{at the right row, the singlet of the three antisymmetrized core qutrit copies} , \end{gathered}$} \end{equation}
and at zero of sixteen at the other assignment. The first parting row the neutral body at its left row, an anti-triplet with the weights $(0, 2)$ derived against the doublet with the weights $(3, -3)$ listed. Over the sixty readings of the grid at both assignments no reading reaches sixteen of sixteen, the best four of sixteen, zero of one hundred twenty pairs complete. The electric charge as the sum of the two weights re-met against the dart census at the dock face at ten of sixteen listed rows, the six up-quark rows the parting rows, and at four of sixteen derived rows. So the sixteen rows are reached at the wide union of the readings and at no one reading, which is the open result. The smallest instance on the page is the electron, its two reads reach both of its rows, one dock entered the left doublet and both docks the right singlet, the weights above, and its color multiplicity one.

The counter-tests struck the open reading of the electron, which moved its doublet cells, and the derivation stood at its five counter-tests and eleven controls. The rows and the comparisons were reproduced at a second engine, equal at every cell.
\chcert{120}
\end{chproof}
\section{The representation census at the saturated sinks, second count}\label{mod:IV-3-D536028-w-rcsy}
\begin{chwall}[The representation census at the saturated sinks, second count]\label{box:IV-3-D536028-w-rcsy:1}
At the derived saturated sinks the open-end census, the kernel counts, the charges, the valence bound, the type closure and the ladder are outputs at every one of the sixteen thousand nine hundred sixty-two sinks, with no listed site as an input, and the list of the board at its names and grade tokens is re-met as that census at four of the six derived entries of the family; the first parting is the generation label entry, where the hub count at the lepton-form sinks outputs two generations, the third label absent at this material, with the quark labels present at three of four. The open result is entered with both cells named, the proton sinks an initial segment at the budget of their entry.
\end{chwall}
\chstanding{open result}
\chdeps{Definitions 1.1 to 1.6 and Open results 66.1 to 66.3, Theorem 66.4 and Open result 66.5}{no theorem beyond it at this release}
The saturated sinks are listed at every id and at every sink the census quantities are computed as outputs; the list of the board is compared with that census entry by entry of the family, re-met at four entries and parting first at the generation label; the two engines agree at every key and the runs are byte-identical.

\begin{chbridge}
The Standard Model lists its fields by their representations under the gauge group, their charges and their generation label, three generations of quarks and leptons with the same gauge quantum numbers repeated at each, the generation count a fact of the list and an output of no principle of the model \cite{Glashow1961,Kobayashi1973}. The derivation supplies the open-end census, the kernel counts and the charges as outputs at every sink here, the gauge group as a corollary standing at chapter 4-2, and the group, the generations, and the charges are named by correspondence.
\end{chbridge}
\begin{chproof}[census]
The premises are the saturation with its sinks (chapter 1-5, a theorem), the representations from the open-end census (chapter 4-3, an open result read at its standing), the conserved functionals as the kernel of the move matrix (chapter 4-6, a theorem) and the list of the board at its names and grade tokens (chapter 4-10, an instrument). The chain back to the axiom runs through chapter 1-5 and the framework chapter. By exhaustive enumeration of the sixteen thousand nine hundred sixty-two saturated sinks of the proton saturation (the set generated by the admitted moves from the seeds to the budget of that derivation, an initial segment and nothing read as a closure). The property tested at every sink being the census quantities, the open-end census, the kernel counts, the charges, the valence bound, the type closure and the ladder, computed as outputs with no listed site as an input, and the list of the board at its names and grade tokens compared with that census entry by entry of the family. The counts by kind, the census an output at every one of the sixteen thousand nine hundred sixty-two sinks and undefined at none, the sink of the list equal to its own block at one of one, the twelve family entries located at twelve of twelve and the nine generation entries at nine of nine, the six derived entries re-met at four, the gauge group as the three algebras located at the end set, the content census of the doublets, the left and right entries reached, and the sixth charge defined at two thousand eight hundred forty-nine of two thousand eight hundred forty-nine sinks carrying a state, and parting at two, the generation label entries of the leptons and the quarks,
\begin{equation}
\text{generation count at the lepton-form sinks by the hub count} = 2,\qquad \text{quark labels present } 3 \text{ of } 4 ,
\end{equation}
the first parting element the lepton generation entry, the labels one and two present at the lepton-form sinks with two hubs at five sinks and one hub at two and the third label absent at this material, which is the open result. The six silent and unmapped entries printed beside at their outputs, among them the anomaly sums over the words at the derivation face, three, one hundred eighty-nine, minus three and zero, and over the sixteen all zero. The smallest case on the page is the sink of the list itself, its census equals its own block, the one sink compared with its listed reading and parting nowhere.

The counter-tests took the census at the listed sites in place of the sinks, which moved fifty-four cells beginning at the gauge-group entry. replaced the kernel by the commutant, which moved twenty-one cells beginning at the baryon entry. cut one link, which moved the first sink, and struck the text at its gate, which left the gate unmet. Every control read the same at both engines, the two engines equal at every key with the runs byte-identical, and the census was reproduced at a second engine, equal at every sink.
\chcert{16962}
\end{chproof}
\section{The representation census at the saturated sinks, third count}\label{mod:IV-3-D538265-w-rcs4}
\begin{chwall}[The representation census at the saturated sinks, third count]\label{box:IV-3-D538265-w-rcs4:1}
At the union of the saturated sinks, the sixteen thousand nine hundred sixty-two sinks of the second pass together with the five thousand three hundred forty tau-entry sinks beyond them, the open-end census, the kernel counts, the charges, the valence bound, the type closure, the ladder and the attach-site count are outputs at every one of the twenty-two thousand three hundred two sinks: the generation count at the lepton-form sinks is three at both readings, by the hub count and by the attach-site count, so the generation parting of the second pass dissolves; the centre weight over six at every open qutrit end lies in the adopted hypercharge canon; but the doublets at the three lepton sinks part at the first, the electron sink with two docks reading the weak singlet by the dock count, and the four anomaly sums at the derivation face are not zero. The open result is entered at the doublets, with both faces named.
\end{chwall}
\chstanding{open result}
\chdeps{Definitions 1.1 to 1.6, Open results 66.1 to 66.3, Theorem 66.4 and Open result 66.5 and Theorems 197.1 to 197.4 and Open result 197.5}{no theorem beyond it at this release}
The union of the lepton-form sinks of the base results is built, the generation count is read at both the hub-count and the attach-site readings, the qutrit ends at every sink are listed and charged, the doublets are read at the three lepton sinks, and the four sums are taken at the derivation face; the counter-tests are met and the controls agree at the pinned grammar, one control re-keyed at a new-name tool by the review at zero freedom; nothing is read as a closure, the base material an initial segment at the budgets of its entries.

\begin{chbridge}
The fermions of the Standard Model fall into three generations, each a copy of the same set of representations of the gauge group, the generation count fixed by nothing within the theory and read from the invisible width of the Z and from the anomaly cancellation within each generation \cite{Glashow1961,Weinberg1967,Bouchiat1972}. The derivation supplies the open-end census, the charges and the four sums as outputs at every sink here, and the doublets, the hypercharge canon and the anomaly are named by correspondence.
\end{chbridge}
\begin{chproof}[census]
The premises are the saturation with its sinks (chapter 1-5, a theorem), the lepton-form sinks of the tau-sink saturation (an open result placed at no chapter of this compendium, read at its standing), the representation census at the saturated sinks at its second count (chapter 4-3, an open result read at its standing) and the Cartan weights at the open ends (chapter 4-4, an open result read at its standing). The chain back to the axiom runs through chapter 1-5 and the framework chapter. By exhaustive enumeration of the twenty-two thousand three hundred two saturated sinks of the union (the set generated by the sixteen thousand nine hundred sixty-two sinks of the second count together with the five thousand three hundred forty tau-entry sinks new beyond them, an initial segment at the budgets of its entries and nothing read as a closure). The property tested at every sink being the census quantities as outputs, the open-end census, the kernel counts, the charges, the valence bound, the type closure, the ladder and the attach-site count, and at the lepton-form sinks the five clauses of the item. The counts by kind, the generation count at the lepton-form sinks three by the hub count and three by the attach-site count, the three lepton labels present at three of three at both readings, so the generation parting of the second count dissolves. The center weight over six at every open qutrit end in the hypercharge canon at eighty-four thousand two hundred fifty-one of eighty-four thousand two hundred fifty-one. The doublets at the three lepton sinks a doublet by the dock count at two of three, and the four anomaly sums at the derivation face
\begin{equation}
(\text{lin},\ \text{cub},\ \text{mix}_{2},\ \text{mix}_{3}) = (3,\ 189,\ -3,\ 0) \ \text{against zero at every one} ,
\end{equation}
the clauses met at three of five. The first parting element in the frozen order the doublets, at the electron seeded sink with two docks reading the weak singlet by the dock count against the doublet of the canon, which is the open result. The derived entries of the board re-met at five of six beside, parting at the quark generation entry alone with the quark labels at three of four. The smallest case on the page is the electron seeded sink, its two docks give the content one, a singlet by the dock count, where the canon asks a doublet.

The counter-tests took the census at the listed sites in place of the sinks, which moved eighty cells. replaced the kernel by the commutant, which moved thirty-three. cut one link, which moved seventy. struck the text at its gate, which left the gate unmet. struck the tau-entry sinks, which moved one hundred twenty-one cells and re-met the parting of the second count at four of four. mistyped a lepton-form sink, which moved the generation count, and struck one word from the sixteen, which moved the census cell. Every control read the same at both engines at the pinned grammar but one, mis-posed at its keying and re-keyed at a new tool at eleven of eleven with no cell moved, and the census was reproduced at a second engine, equal at every sink.
\chcert{22302}
\end{chproof}
\section{The carrier sink offices}\label{mod:IV-3-D540063-cbsk-t1}
\begin{chtheorem}[The carrier sink offices]\label{box:IV-3-D540063-cbsk-t1:1}
The charged, neutral and Higgs-form sinks of the saturation census beyond the lepton and nucleon forms are bodies at their own sinks under the one rule, their offices exact cells in the electron unit at every hub count, with the carrier words of the Higgs and carrier row beside them as the word faces: the census of seventeen thousand seven hundred three sinks sorts by form into sixteen thousand seven hundred eighty-eight quark, seven hundred forty-seven lepton and one hundred sixty-eight dark, and by type into one thousand twenty charged, two thousand two hundred two neutral, one hundred sixty-eight dark and fourteen thousand three hundred thirteen beyond two docks, the attachment standing at ten thousand one hundred twenty-three; at the adopted reading the material is the two thousand eight hundred forty-two sinks of the proton-web saturation beyond the lepton and nucleon forms, seven hundred eighty-nine charged, one thousand eight hundred eighty-six neutral and one hundred sixty-seven dark, every one located at its anatomy and typed by its dock count at two thousand eight hundred forty-two of two thousand eight hundred forty-two, and the office under the one rule at the measure-struck reading is defined as an exact nonzero element of the field at two thousand eight hundred forty-two of two thousand eight hundred forty-two, every one with a root part, the first power from seven hundred forty-seven to twenty-four million two hundred ten thousand ninety; the first charged sink, at hub count one, twelve sites, fourteen links, the measure one hundred thirteen, one dock struck and the dressing twenty-seven, has the first power two thousand eight hundred eighty-nine and the office $\tfrac{5\ 469\ 858\ 260\ 402\ 743}{655360000} + \tfrac{649}{50}\sqrt{2} + \tfrac{6626373}{32768000}\sqrt{3}$, the first neutral sink the first power two thousand five hundred sixty-five and the first dark sink three thousand two hundred thirteen; the Higgs-form sinks at ten thousand one hundred twenty-three are defined at ten thousand one hundred twenty-three of ten thousand one hundred twenty-three, the first within the material at the first power three thousand three hundred seventy-five and the office $\tfrac{186\ 625\ 265\ 011\ 095\ 391}{16384000000} + \tfrac{1268}{25}\sqrt{2} + \tfrac{6626373}{32768000}\sqrt{3}$; the word faces re-met at the Higgs and carrier row are the W and Z words at the first power twenty-eight and the office $\tfrac{2857690241}{3645000} + \tfrac{98}{50625}\sqrt{2}$ each and the Higgs word at the first power sixteen and the office $\tfrac{46656001681}{182250000}$, every office over the electron office an exact element compared with nothing. The first face reads yes at the adopted reading and the second face, the first sink the rule refuses named, reads no, the first parting none.
\end{chtheorem}
\chstanding{theorem}
\chdeps{Definitions 1.1 to 1.6, Open results 66.1 to 66.3, Theorem 66.4 and Open result 66.5, Theorems 185.1 to 185.12, Theorem 195.1, Theorems 200.1 to 200.3 and Open result 200.4 and Open results 245.1 and 245.2, Theorem 245.3, Open results 245.4 to 245.8 and Theorem 245.9}{no theorem beyond it at this release}
The census is re-run from the copies and sorted by form and by type; the material is located at its anatomy in the webs table; the office of every sink is read under the one rule at the measure-struck reading with the dock set struck, the dark set every link beyond the (2,3) genre and the closure at a dark link the square of its act multiplied by the second-order read, and ordered by hub count and form; the Higgs-form sinks are read alike with the vacuum object beside, and the word faces are re-met at the Higgs and carrier row.

\begin{chbridge}
The masses of the weak bosons and of the Higgs boson arise from the vacuum expectation value of the Higgs doublet, while the hadrons take their masses from the binding of the strong force, so that the carriers and the bound states are placed on one scale in the unit of the electron \cite{Weinberg1967,GellMann1964}. The derivation supplies the offices of the carrier sinks and the word faces here, and the masses are named by correspondence.
\end{chbridge}
\begin{chproof}[census and exact]
The premises are the terminated saturation at the census (chapter 1-5, a theorem), the body-alone office rule at the measure-struck reading (chapter 8-2, a theorem), the Higgs and carrier offices with the word faces (chapter 6-10, a theorem), the Higgs attachment at the (4,4) count (chapter 6-10, an open result) and the count law at the six genres (chapter 8-1, a theorem). The chain back to the axiom runs through chapter 1-5 and the framework chapter. By exhaustive enumeration of the seventeen thousand seven hundred three sinks of the census, sixteen thousand nine hundred forty-six of the proton-web saturation, seven hundred thirty-eight of the tau sink, fourteen of the roster, three of the recorded webs and two beside, the property tested at each being whether the one rule defines its office as an exact nonzero element and whether its type by the dock count of the base agrees with its anatomy. From the material at the adopted reading, the two thousand eight hundred forty-two sinks beyond the lepton and nucleon forms, by the type at the dock count, we obtain the charged at seven hundred eighty-nine at the hub counts one to six, the neutral at one thousand eight hundred eighty-six at one to eight and the dark at one hundred sixty-seven at one to six, by wave five, one hundred one, two thousand forty-four and six hundred ninety-two, the three groups present and the Higgs form present at six hundred sixty-five within the material and nine thousand four hundred fifty-eight beyond. From the first charged sink, by the one rule, the act of a link its measure multiplied by the hub dressing at every unstruck link, the dock set struck, the dark set every link of a genre other than (2,3) and the closure at a dark link the square of its act multiplied by the second-order read, we obtain
\begin{equation}
p = 2889 , \qquad p^{2} = 8346321 , \qquad c = \tfrac{13329842743}{655360000} + \tfrac{649}{50}\sqrt{2} + \tfrac{6626373}{32768000}\sqrt{3} ,
\end{equation}
the closure positive at nine dark links and one hundred seventy-seven closure records, and the office $p^{2} + c = \tfrac{5\ 469\ 858\ 260\ 402\ 743}{655360000} + \tfrac{649}{50}\sqrt{2} + \tfrac{6626373}{32768000}\sqrt{3}$; alike the first neutral sink at the first power two thousand five hundred sixty-five and the office $\tfrac{4\ 311\ 774\ 225\ 842\ 743}{655360000} + \tfrac{649}{50}\sqrt{2} + \tfrac{6626373}{32768000}\sqrt{3}$ and the first dark sink at three thousand two hundred thirteen and $\tfrac{6\ 765\ 536\ 437\ 682\ 743}{655360000} + \tfrac{649}{50}\sqrt{2} + \tfrac{6626373}{32768000}\sqrt{3}$, the twenty groups by hub count and form each with its first sink. The counts by kind: the office defined at two thousand eight hundred forty-two of two thousand eight hundred forty-two with the two-read identity at every read; the Higgs form defined at ten thousand one hundred twenty-three of ten thousand one hundred twenty-three, the first power from three hundred thirty-six to eight hundred sixty-eight million nine hundred six thousand thirty-five beyond the material, the quartic spelling of every Higgs-form sink, its office over the vacuum object squared, exact at every one; the word faces by the same rule re-met at the Higgs and carrier row at three of three, the W and Z words at the first power twenty-eight and the Higgs word at sixteen; the sum of the first powers over the material and the Higgs form beyond three hundred forty-two billion six hundred sixty-nine million two hundred sixty-six thousand six hundred seventy. So the carrier types are bodies at their own sinks with their offices exact cells, which is the theorem. The smallest instance on the page is the Higgs word: one (4,4) link at the first power sixteen, the office $\tfrac{46656001681}{182250000}$.

The counter-tests left the dock set unstruck, which moved the sum of the first powers and the sum of the offices. A second read the dark set as every link, which moved the sum of the closures, the word faces and the deciding cells with the word faces parting from the Higgs and carrier row. A third struck the attachment from the dark set, which moved the closures and the Higgs word face, a fourth struck the dressing to one, which moved the first powers, and a fifth struck the gate, which left every cell absent. Every control read the same at both engines, thirteen of thirteen, and the offices were reproduced at a second engine of a different construction, byte-identical at every run.
\chcert{12300}
\end{chproof}
\begin{figure}[h]\centering
\begin{adjustbox}{max width=\linewidth}\begin{tabular}{lrrrrrr}\toprule
the first sink of its kind & hub count & sites & links & measure & dressing & first power \\ \midrule
charged & 1 & 12 & 14 & 113 & 27 & 2889 \\
neutral & 1 & 11 & 13 & 107 & 27 & 2565 \\
dark & 1 & 13 & 15 & 119 & 27 & 3213 \\
Higgs form within the material & 2 & & & & 27 & 3375 \\ \bottomrule
\end{tabular}\end{adjustbox}
\caption{The first sink of each kind in the census order at the adopted reading, with its hub count, its sites and links, its measure, its dressing and its first power.}
\end{figure}
\section{The same power at the shifted origin}\label{mod:IV-3-D543926-foso-t1}
\begin{chtheorem}[The same power at the shifted origin]\label{box:IV-3-D543926-foso-t1:1}
The carrier chains of the Green-function row fall by the same first power as the pure chains once the origin is shifted by the (2,3) register of the recorded link: the printed elements of that row at every separation to twelve at seven chains are, at the two carrier chains, the proton sink at its (2,3) link and the photon word at its (2,3) link, $\tfrac{79 + 48\sqrt{2} + 42\sqrt{3} + 3\sqrt{6}}{31104}$, over thirty-three thousand four hundred eight, over thirty-five thousand seven hundred twelve and over thirty-eight thousand sixteen at the separations one to four and over forty thousand three hundred twenty to over fifty-six thousand four hundred forty-eight at five to twelve, and at the five pure chains, the electron, the muon, the tau, the deuteron and the W word at a (2,2) link, $\tfrac{30 + 30\sqrt{2}}{2304}$ multiplied by the separation; the ratios of successive separations recomputed as exact elements of the ring are $\tfrac{27}{29}, \tfrac{29}{31}, \tfrac{31}{33}$ and on to $\tfrac{47}{49}$ at the carrier chains and $\tfrac{1}{2}, \tfrac{2}{3}$ and on to $\tfrac{11}{12}$ at the pure chains, equal to the printed ratios at every printed separation; the (2,3) register of the recorded link is thirty-one thousand one hundred four at the Ward-functional row, at the charge register of the Green-function row and at the first step of both carrier chains, the (2,2) register of the carrier two thousand three hundred four, their ratio $\tfrac{27}{2}$ and the origin offset $\tfrac{25}{2}$ at the separation convention counting the recorded link as one link, so that the register of the carrier chain at the separation $k$ is two thousand three hundred four multiplied by $k + \tfrac{25}{2}$; the first-power law at the shifted origin, the ratio at the separation $k$ the quotient of $k - 1$ plus the offset over $k$ plus the offset, equals the recomputed ratios at eleven of eleven separations at both carrier chains at the offset $\tfrac{25}{2}$ and at the five pure chains at the offset zero, seven of seven chains equal at seventy-seven of seventy-seven separations. The first face reads yes and the second face, the first separation parting named, reads no, the first parting none.
\end{chtheorem}
\chstanding{theorem}
\chdeps{Definitions 1.1 to 1.6, Open results 66.1 to 66.3, Theorem 66.4 and Open result 66.5, Theorems 78.1 to 78.3, Open result 92.1 and Theorem 92.2 and Theorem 115.1}{no theorem beyond it at this release}
The printed elements of the Green-function row are located at every separation to twelve at the seven chains and their successive ratios are recomputed as exact elements of the ring; the (2,3) register of the recorded link is re-met at three locators and the origin offset is read at the separation convention; the first-power law at the shifted origin is evaluated at every separation and compared with the ratios at every chain.

\begin{chbridge}
A propagator on a line falls as the inverse of the separation, and a source attached at the origin shifts the separation by a constant offset without changing the power, the Green function of a point source \cite{Feynman1949,Schwinger1951}. The derivation supplies the ratios at every separation and the offset here, and the propagator is named by correspondence.
\end{chbridge}
\begin{chproof}[finite and exact]
The premises are the Green-function row at its two-point table and its charge table (chapter 10-2, an open result), the two-two two-point function at its register (chapter 10-2, an open result), the Ward-functional row at the (2,3) register (chapter 4-8, an open result) and the serial insertion row at the chains (chapter 10-6, an open result). The chain back to the axiom runs through chapter 10-0 and the framework chapter. The procedure. The input is the seven chains, the proton sink and the photon word at their (2,3) links and the five pure chains at a (2,2) link, with the printed elements of the Green-function row at the separations one to twelve. The enumeration rule locates every printed element, recomputes the ratio of successive separations as an exact element of the ring, re-meets the (2,3) register at its three locators, reads the offset at the separation convention and evaluates the first-power law at the shifted origin at every separation; the test at each chain and separation is equality of the law with the recomputed ratio. From the registers of the recorded (2,3) link and of the (2,2) carrier, by the quotient, we obtain the offset
\begin{equation}
\frac{31104}{2304} = \frac{27}{2} , \qquad \delta = \frac{27}{2} - 1 = \frac{25}{2} , \qquad R_{k} = 2304 \left( k + \tfrac{25}{2} \right) ,
\end{equation}
the register law of the Green-function row re-met at fifty-two of fifty-two; and from the law at the shifted origin, by the ratio,
\begin{equation}
\frac{G_{k}}{G_{k-1}} = \frac{k - 1 + \delta}{k + \delta} = \frac{27}{29}, \frac{29}{31}, \ldots, \frac{47}{49} \quad \text{at } k = 2, \ldots, 12 ,
\end{equation}
equal to the recomputed ratios at eleven of eleven at the proton chain and at the photon word, and at the offset zero $\tfrac{1}{2}, \tfrac{2}{3}, \ldots, \tfrac{11}{12}$ at the five pure chains. The counts by kind: every element exact at twelve of twelve separations at seven of seven chains, the first not exact none; the printed ratios equal to the recomputed at three of three at each carrier chain and at three of three and eleven of eleven at the pure chains; the (2,3) register equal at every locator with the offset exact, the inserted-link convention at $\tfrac{27}{2}$ printed beside; the element law beside, the numerator over two thousand three hundred four multiplied by $k$ plus the offset, equal to the printed element at twelve of twelve at every chain; the chains equal at seventy-seven of seventy-seven separations. So the carriers carry the same law as the pure chains and the (2,3) charge propagates by the same count, which is the theorem. The smallest instance on the page is the proton chain at the second separation: the ratio $\tfrac{27}{29}$, the law at $\tfrac{1 + 25/2}{2 + 25/2}$.

The counter-tests struck the offset to zero at the carrier chains, which moved the law off the ratios at the first separation. A second read the offset at the inserted-link convention, $\tfrac{27}{2}$, which moved the first parting from none to the second separation. A third read a second-power law, which moved every chain off the ratios, a fourth replaced the printed proton ratio at the second separation by the one at the third, which moved the equality to no, and a fifth struck the gate, which left every cell struck. Every control read the same at both engines, twelve of twelve, and the ratios were reproduced at a second engine of a different construction, byte-identical at every run.
\chcert{77}
\end{chproof}
\begin{figure}[h]\centering
\begin{adjustbox}{max width=\linewidth}\begin{tabular}{lllll}\toprule
chain & the offset & the element at the separation $k$ & the ratio at $k = 2$ & the ratio at $k = 12$ \\ \midrule
the proton sink at its (2,3) link & $\tfrac{25}{2}$ & $\tfrac{79 + 48\sqrt{2} + 42\sqrt{3} + 3\sqrt{6}}{2304\,(k + 25/2)}$ & $\tfrac{27}{29}$ & $\tfrac{47}{49}$ \\
the photon word at its (2,3) link & $\tfrac{25}{2}$ & the same element & $\tfrac{27}{29}$ & $\tfrac{47}{49}$ \\
the electron at its (2,2) link & $0$ & $\tfrac{30 + 30\sqrt{2}}{2304}\, k$ & $\tfrac{1}{2}$ & $\tfrac{11}{12}$ \\
the W word at its (2,2) link & $0$ & $\tfrac{30 + 30\sqrt{2}}{2304}\, k$ & $\tfrac{1}{2}$ & $\tfrac{11}{12}$ \\ \bottomrule
\end{tabular}\end{adjustbox}
\caption{The element of the Green-function row at the separation $k$ and its successive ratios at four of the seven chains, the first-power law at the shifted origin at the carriers and at the zero offset at the pure chains.}
\end{figure}
\section{The wave as the carrier advance and the removing moves}\label{mod:IV-3-D543954-the-wave-as-the-carrier-advance-and-the-removing-moves}
\begin{chtable}[The wave as the carrier advance and the removing moves]\label{box:IV-3-D543954-the-wave-as-the-carrier-advance-and-the-removing-moves:1}
The (2,3) modulation along the rung ladders and the removing moves of the horizon record are printed rung by rung as two series side by side: the three rung ladders of the wave row, two at the base of fifteen links with three (2,2) and one (2,3) link and the third at the universe web of five hundred fifty-one links with eighty-five (2,2) and thirty-four (2,3) links, carry a (2,3) link at every one of their six rung links, the (2,3) count rising one, two, three, four, five, six, seven along the rungs zero to six at the first two ladders and thirty-four to forty at the third with the difference one at every step, uniform, while the (2,2) count stands still at three and at eighty-five, re-met at the wave row at thirty-six of thirty-six cells; the horizon record, every recorded web parsed from its state line with its membership vector, holds the webs by wave one, seven hundred ninety-two and thirty-seven thousand two hundred seventy-four at the adopted reading and one, seven hundred thirty-eight and thirty-five thousand sixteen at the reading beside, the tick along the parent column equal to the printed wave at thirty-eight thousand sixty-seven of thirty-eight thousand sixty-seven; the step of every web against its parent is typed, and at the first rung the (2,2) links removed are six hundred seventy-eight at six hundred seventy-one steps with eight added, the net minus six hundred seventy and the count of the webs from eighty-three to eighty-six, at the second rung thirty-two thousand eight hundred seventy-three removed at thirty-two thousand five hundred fifty-seven steps with twenty-three added, re-met at the abundances row at eighty-nine of eighty-nine step kinds, four of four growths and six of six rungs; the rays from the seed along the parent column to every web of the second rung, thirty-seven thousand two hundred seventy-four at the adopted reading and thirty-five thousand sixteen beside, each of two differences, are periodic at the adopted rule, a non-constant pattern repeating within the ray, at none, so that the alternating pattern is not witnessable at rays of two differences, with the constant words at three thousand eight hundred five and the adjacent differences of opposite signs at five thousand five hundred printed beside; the order sequences per rung are aperiodic at both rungs at both readings, the distinct differences minus two, minus one, zero and one. The table is printed as an instrument of the chapter.
\end{chtable}
\chstanding{instrument}
\chdeps{Definitions 1.1 to 1.6, Theorems 9.1 to 9.7 and Open result 9.8, Open results 66.1 to 66.3, Theorem 66.4 and Open result 66.5, Theorems 113.1 and 113.2 and Open results 245.1 and 245.2, Theorem 245.3, Open results 245.4 to 245.8 and Theorem 245.9}{no theorem beyond it at this release}
The three rung ladders are read at their (2,3) and (2,2) counts rung by rung and re-met at the wave row; the horizon record is parsed web by web with the step of every web typed against its parent and the removed and added (2,2) links counted at the membership vectors; the two series are printed side by side; every ray of two differences is read for its word at the adopted rule and at the readings beside, and the order sequences are read per rung.

\begin{chbridge}
A wave carries its disturbance forward one site per step while the medium behind the front returns toward rest, so that the advance of the front and the recoil of the medium are two series read side by side \cite{Brillouin1946,Feynman1963}. The derivation supplies the two series at every rung here, and the wave is named by correspondence.
\end{chbridge}
\begin{chproof}[instrument]
The premises are the waves at the two-two web with its ladders (chapter 1-5, an open result), the termination census as the abundances (chapter 1-5, a theorem), the horizon depth at its record (chapter 10-7, an open result) and the tensor type at every genre with its increment (chapter 1-2, an open result). The chain back to the axiom runs through chapter 1-5 and the framework chapter. The procedure. The input is the three rung ladders of the wave row, the horizon record at its state lines and the tensor-type increment. The enumeration rule reads the (2,3) and (2,2) counts along every ladder, parses every recorded web with its membership vector and types its step against its parent, counts the (2,2) links removed and added per rung at the membership vectors since the canonical relabelling of the recorded webs scrambles the site labels at some children, prints the two series side by side, reads every ray for its minimal repeating word and reads the order sequences per rung; no cell is computed beyond a count. The counts by kind: the ladders three with the rung links (2,3) at six of six each and the cells re-met thirty-six of thirty-six; the webs thirty-eight thousand sixty-seven at the adopted reading and thirty-five thousand seven hundred fifty-five beside with the tick equal to the wave at every one; the step kinds at the first rung the join at (2,2) eight, the reversed attach one, the reversed cut six hundred thirty, the reversed hub step at its three kinds, the reversed hub forms and the reversed join at (3,3) fifty-four; the rays by word at the adopted reading seven words with the constant words three thousand eight hundred five and the sign-alternating words five thousand five hundred; the periodic rays none of thirty-seven thousand two hundred seventy-four and none of thirty-five thousand sixteen. The smallest instance on the page is the first ladder at its first step: the (2,3) count from one to two, the (2,2) count three still, the links from fifteen to sixteen.

The counter-tests read the rays and the order sequences backward, which moved thirteen of three hundred fifty-seven shared cells with the period letters standing. A second struck the reversed cut kind, which moved one hundred eight cells with the removals and the ray words. A third put the (2,3) count in place of the (2,2) along the rays, which moved ninety-seven cells, and a fourth struck the gate, which left no cell built at both engines. Every control read the same at both engines, six of six, and the series were reproduced at a second engine, equal at every cell.
\end{chproof}
\par\medskip\noindent{\itshape From the representations the thread passes to their weights and the charge.}\label{narr:close:4-3}\par
\section{The Cartan weights and the charge}\label{mod:IV-4-cartan-weights-and-the-charge}
\chdesig{5-4}{4-4}{IV-4-cartan-weights-and-the-charge}{IV-4-cartan-weights-and-the-charge}
\label{narr:open:4-4}Open result~72.1 reads the Cartan weights and the charge, Open result~72.2 the representation table at the seeded sinks, Open result~72.3 the charge as a centralizer weight at the dock face, Theorem~72.4 the signed open-end count under the persistence matching, and Open result~72.5 the Cartan functionals and the dart at the second-order face, the conditional result among open ones.
\begin{chwall}[The Cartan weights and the charge]\label{box:IV-4-cartan-weights-and-the-charge:1}
This chapter assigns the Cartan weights (the Cartan weight) of the construction at every open end and finds the charge a centraliser weight at the dock face, while the electric charge of the dart (the dart: a directed edge of a graph) is returned as an output at not every web. The weights assigned at every open end of every listed web and every sink, the weak weight at the hub-leg qubit copies by the end letters of the network list, the strong weights at the qutrit ends and the phase weight at the dock ends, meet the dart at four of the twenty-two listed webs carrying an open end and part first at the docked up web, while the sixths combinations hold at every web, the four trace identities vanish over the sixteen rows, and the sign at the muon agrees.
\end{chwall}
\chstanding{open result}
\chdeps{Definitions 1.1 to 1.6, Theorems 9.1 to 9.7 and Open result 9.8, Theorems 64.1 and 64.2, Theorems 65.1 and 65.2, Open results 66.1 to 66.3, Theorem 66.4 and Open result 66.5, Theorem 163.1, Open result 163.2 and Theorem 163.3, Open results 241.1 and 241.2, Theorems 241.3 and 241.4 and Open results 241.5 to 241.7 and Open results 245.1 and 245.2, Theorem 245.3, Open results 245.4 to 245.8 and Theorem 245.9}{Theorem 73.1 and Open results 73.2 and 73.3, Theorems 74.1 to 74.4 and Open result 74.5, Open results 76.1 to 76.3, Theorems 78.1 to 78.3 and Open results 84.1 to 84.3}
\begin{chwall}[The representation table at the seeded sinks]\label{box:IV-4-cartan-weights-and-the-charge:2}
At the seeded sinks the representation table of the construction is printed at every sink and at the released object of the decay graph, which reads as the lower member of the doublet at hypercharge zero under both rules named and as the electron row under the pendant reading beside, the owed word returned at no row.
\end{chwall}
\chstanding{open result}
\chdeps{Definitions 1.1 to 1.6, Theorems 9.1 to 9.7 and Open result 9.8, Theorems 64.1 and 64.2, Theorems 65.1 and 65.2, Open results 66.1 to 66.3, Theorem 66.4 and Open result 66.5, Theorem 163.1, Open result 163.2 and Theorem 163.3, Open results 241.1 and 241.2, Theorems 241.3 and 241.4 and Open results 241.5 to 241.7 and Open results 245.1 and 245.2, Theorem 245.3, Open results 245.4 to 245.8 and Theorem 245.9}{Theorem 73.1 and Open results 73.2 and 73.3, Theorems 74.1 to 74.4 and Open result 74.5, Open results 76.1 to 76.3, Theorems 78.1 to 78.3 and Open results 84.1 to 84.3}
\begin{chwall}[The charge as a centralizer weight at the dock face]\label{box:IV-4-cartan-weights-and-the-charge:3}
The charge lies inside the centraliser of the dock-face stabiliser and outside the centraliser of the construction whole, so it is a centraliser weight at the dock face, and the dart re-read in the signed-dart-end basis equals that weight at seventeen of the twenty-two listed webs at the dock-rule face and at fifteen at the hub-leg face.
\end{chwall}
\chstanding{open result}
\chdeps{Definitions 1.1 to 1.6, Theorems 9.1 to 9.7 and Open result 9.8, Theorems 64.1 and 64.2, Theorems 65.1 and 65.2, Open results 66.1 to 66.3, Theorem 66.4 and Open result 66.5, Theorem 163.1, Open result 163.2 and Theorem 163.3, Open results 241.1 and 241.2, Theorems 241.3 and 241.4 and Open results 241.5 to 241.7 and Open results 245.1 and 245.2, Theorem 245.3, Open results 245.4 to 245.8 and Theorem 245.9}{Theorem 73.1 and Open results 73.2 and 73.3, Theorems 74.1 to 74.4 and Open result 74.5, Open results 76.1 to 76.3, Theorems 78.1 to 78.3 and Open results 84.1 to 84.3}
\begin{chconditional}[The signed open-end count under the persistence matching]\label{box:IV-4-cartan-weights-and-the-charge:4}
Conditional on the one-act clause, and on the results these theorems rest on at their own standing.
Under the persistence matching the signed open-end count at the unit weight is conserved at all but eight hundred thirty-five of the move instances of the census, every changing instance a re-matching of two ends of one type, so the conservation is a theorem of the joins and an open result of the re-matchings, conditional on the one-act clause.
\end{chconditional}
\chstanding{conditional theorem}
\chdeps{Definitions 1.1 to 1.6, Theorems 9.1 to 9.7 and Open result 9.8, Theorems 64.1 and 64.2, Theorems 65.1 and 65.2, Open results 66.1 to 66.3, Theorem 66.4 and Open result 66.5, Theorem 163.1, Open result 163.2 and Theorem 163.3, Open results 241.1 and 241.2, Theorems 241.3 and 241.4 and Open results 241.5 to 241.7 and Open results 245.1 and 245.2, Theorem 245.3, Open results 245.4 to 245.8 and Theorem 245.9}{Theorem 73.1 and Open results 73.2 and 73.3, Theorems 74.1 to 74.4 and Open result 74.5, Open results 76.1 to 76.3, Theorems 78.1 to 78.3 and Open results 84.1 to 84.3}
\begin{chwall}[The Cartan functionals and the dart at the second-order face]\label{box:IV-4-cartan-weights-and-the-charge:5}
The count of Cartan functionals conserved at every census entry equals the Cartan dimension at three of the six link types; and the up body carries its electric dart at the second-order face while the sixteen rows read equal at no face, the first parting moving from the first up row to the first down row.
\end{chwall}
\chstanding{open result}
\chdeps{Definitions 1.1 to 1.6, Theorems 9.1 to 9.7 and Open result 9.8, Theorems 64.1 and 64.2, Theorems 65.1 and 65.2, Open results 66.1 to 66.3, Theorem 66.4 and Open result 66.5, Theorem 163.1, Open result 163.2 and Theorem 163.3, Open results 241.1 and 241.2, Theorems 241.3 and 241.4 and Open results 241.5 to 241.7 and Open results 245.1 and 245.2, Theorem 245.3, Open results 245.4 to 245.8 and Theorem 245.9}{Theorem 73.1 and Open results 73.2 and 73.3, Theorems 74.1 to 74.4 and Open result 74.5, Open results 76.1 to 76.3, Theorems 78.1 to 78.3 and Open results 84.1 to 84.3}
The unit weight is assigned at every open end of every listed web under the adopted typing, the signed count is summed, and its change is tabulated at every one of the thirty-two thousand eight hundred and fifty-one census entries under the projection matching, the changing instances typed by move kind; the weights of the construction are assigned end by end and compared with the dart and with the sixteen rows; the centralizer of the dock-face stabilizer is computed as a commutant and the charge placed inside or outside it.

\begin{chbridge}
A Cartan subalgebra of a semisimple Lie algebra is a maximal abelian subalgebra of diagonalisable elements, and a representation decomposes into weight spaces labelled by the eigenvalues of its generators, the weights of the fundamental of $su(2)$ being $\pm 1/2$ and those of the fundamental of $su(3)$ the three weights of the triplet \cite{Humphreys1972}; the electric charge of the Standard Model is $Q = T_3 + Y$, the hypercharges of the quarks at the sixths lattice \cite{GellMannNishijima1953,Weinberg1967}. The derivation supplies the weights of the construction at every open end and the signed open-end count here, and the Cartan weights, the electric charge, and the hypercharge are named by correspondence, the conservation standing at chapters 4-6 and 4-7 by the kernel alone.
\end{chbridge}
\begin{chproof}[census]
The premises are the gauge group as a corollary of one construction with its three factors (chapter 4-2, a theorem), the listed webs with their end letters (chapter 1-5, a theorem), the saturation with its sinks (chapter 1-5, a theorem) and the sign from the side of the acting site (chapter 4-5, a theorem). The chain back to the axiom runs through chapter 4-2 and the framework chapter. By exhaustive enumeration of every open end of every listed web and of every one of the forty-four thousand three hundred forty-one sinks of the saturation (the set generated by the admitted moves from the seeds at the budget of that derivation, the site, link and hub counts re-met at every sink). The property tested being whether the five outputs of the weights assigned at the open ends, the sixths combinations, the four trace identities over the sixteen rows, the sign at the muon, the centralizer direction and the charge against the dart, hold at every web. The counts by kind, the weak weight at six hundred fifty-seven hub-leg qubit copies by the end letters of the list, lettered at six hundred fifty-seven of six hundred fifty-seven, two hundred fifteen joined pairs canceling and two hundred twenty-seven odd legs remaining, the strong weights at two hundred nineteen qutrit ends at both sign faces, the phase weight at twenty-five dock ends as the slice entry over six, and at the sinks the orientation letters at five thousand seven hundred seventy-two of five thousand seven hundred seventy-two instances with five thousand one hundred three resolved. The sixths integer combinations at every web, the four trace identities zero over the sixteen rows, the sign at the muon equal and the centralizer direction existing at the quotient dimension one, four outputs of five met. The charge as the sum of the weights equal to the dart at four of the twenty-two listed webs carrying an open end at the dock-rule face and at four of twenty-two at the hub-leg face, at zero of twenty-two at the plus faces. The first parting element in the list order the docked up web, its ends at the hub legs ket plus three, bra minus three and ket plus three and its dock minus three,
\begin{equation}
6Q = 3 - 3 + 3 - 3 = 0 \quad\text{against the dart}\quad -6 \ \text{in sixths} ,
\end{equation}
so the dart alone parts, which is the open result. The smallest case on the page is that docked up web, its four open ends carry the weights above, their sum vanishes, and the dart unit composition reads minus six.

The counter-tests read every hub-leg copy as ket, which moved every weak weight to plus one half and the sign, the charge and the muon cells. set the slice entries to zero at every open end, which moved the hypercharge to zero everywhere and the charge and the covering cells, and struck the text at its gate, which left the gate unmet. Every control read the same at both engines, among them the list parsed in its block order with no line dropped and the dotted dock copies read, and the weights were reproduced at a second engine, equal at every end.
\chcert{44341}
\end{chproof}
\par\medskip\noindent{\itshape The weights carry a sign, and the next chapter says where the sign comes from.}\label{narr:close:4-4}\par
\section{The sign from the side of the acting site}\label{mod:IV-5-sign-from-the-side-of-the-acting-site}
\chdesig{5-5}{4-5}{IV-5-sign-from-the-side-of-the-acting-site}{IV-5-sign-from-the-side-of-the-acting-site}
\label{narr:open:4-5}Theorem~73.1 obtains the sign from the side at the hub reached through its qutrit, with the slot rule from the sink orientation, Open result~73.2, and the slot rule from the exchange-to-end map, Open result~73.3, open beside it.
\begin{chtheorem}[The sign from the side at the hub reached through its qutrit]\label{box:IV-5-sign-from-the-side-of-the-acting-site:1}
This chapter establishes that the isospin sign is the side of the acting site (the acting site) at every hub reached through its qutrit, and that the slot rule is not the sink orientation read at the cut end. Under the side carry from the probe dock, every hub the carry enters through its qutrit leg carries the far side at its odd-leg qubit copy, equal to the letter of the network list there, at all five hubs of the ladder so entered and at every hub of the three lepton-seeded sinks; the tau's second hub is entered through its qubit leg, its odd-leg copy reading the near side as a corollary of its leaf attachment, the one hub outside the domain. The open result that the sign is one function of the side at every hub was entered at that hub and dissolves at the re-wording.
\end{chtheorem}
\chstanding{theorem}
\chdeps{the theorems of Chapters 1-0, 1-2 to 1-6, 5-2, 5-4, 8-3}{Open results 84.1 to 84.3}
\begin{chwall}[The slot rule from the sink orientation]\label{box:IV-5-sign-from-the-side-of-the-acting-site:2}
The adopted slot rule, the end of larger degree, then of larger dimension, then the first-named end, is not the sink orientation read at the cut end under any of the six readings, the first parting the cut at the qutrit of the hub leg of the docked down body; the first cut at which the slot rule and the orientation from the apparatus acting site part is the hub leaf leg of the electron body; a cut of the invariant part under the orientation freedom parts the slot rule, the first the core link of the electron; and under the hub-leg reversal the two slot-rule open results dissolve while a cut of the invariant part still parts the rule beyond the naming clause, the first the attach tether of the electron.
\end{chwall}
\chstanding{open result}
\chdeps{the theorems of Chapters 1-0, 1-2 to 1-6, 5-2, 5-4, 8-3}{Open results 84.1 to 84.3}
\begin{chwall}[Slot rule from the exchange to end map]\label{box:IV-5-sign-from-the-side-of-the-acting-site:3}
At the link, the listed slot end against the centre end of the exchange-to-end map agrees at most lettered links, differs at fifty-one by the degree clause and is undetermined at the seven hundred sixty-one equal-dimension links; the admitted pairing of the two reads with the two exchange ends holds at forty-eight of the one hundred seventy-two eigenvectors; the one-read-per-end count meets the slot-pair clause at three of the nine differing eigenvectors; the polarity of the read line fixes a centre at every free link that equals the derived slot end at about half the lettered links; the far letter pair of the six even eigenvectors equals the ket-and-bra reads at three; the slot end at equal dimension is a washed label at the dimension-seen readouts and moves at nineteen links at the star face; the larger-star end equals the listed slot at thirteen of those nineteen; and the sign at the unread components differs from the relative component sign at the odd set.
\end{chwall}
\chstanding{open result}
\chdeps{the theorems of Chapters 1-0, 1-2 to 1-6, 5-2, 5-4, 8-3}{Open results 84.1 to 84.3}
The carried arrow from the probe dock is followed into every hub of the ladder and of the three lepton-seeded sinks, each hub typed by the leg through which the carry enters it and the side of its odd-leg copy compared with the letter of the network list; the slot rule is applied at every cut of the sixteen words of the act and of the five joined webs and compared with the carried arrow at every reading; the freedoms of the orientation are located and the invariant part isolated; at the link the slot end is compared with the exchange-to-end map, the polarity center and the star face link by link.

\begin{chbridge}
In the electroweak theory the two members of a weak-isospin doublet are distinguished by the sign of $T_3$ \cite{Weinberg1967,StreaterWightman1964}; in a quantum field theory the orientation of a propagator line between a ket index and a bra index fixes the sign of a two-point function, a convention carried consistently through every diagram \cite{Peskin1995}. The derivation supplies the side of the acting site and the weight letters of the network list here, the representations standing at chapter 4-3 by the open-end census alone, and the isospin and its sign are named by correspondence.
\end{chbridge}
\begin{chproof}[exact]
The premises are the probe electron at every measurement web (chapter 1-6, a theorem), the ladder webs of the three leptons and their seeded sinks (chapter 1-5, a theorem), the sink orientation up to the junction reversal (chapter 1-4, a theorem) and the Cartan weights at the open ends (chapter 4-4, an open result read at its standing). The chain back to the axiom runs through chapter 1-6 and the framework chapter. From the probe dock of each lepton web, by the carried arrow followed unreversed from the dock into the web, we obtain at every hub the leg through which the carry enters it and the side of its odd-leg qubit copy, the far side where the carry enters through the qutrit leg and the near side where it enters through the qubit leg,
\begin{equation}
s(h) = \begin{cases} -1 & \text{the carry enters } h \text{ through its qutrit leg (the far side, bra)}\\ +1 & \text{the carry enters } h \text{ through its qubit leg (the near side, ket)} \end{cases}
\end{equation}
and from the network list, by the end letter at the odd-leg qubit copy of every hub, the weight sign $w(h) = -1$ at every one. From the six hubs of the ladder, by the comparison of the two signs hub by hub, we obtain equality at five, the electron hub, the two hubs of the muon and the first and third hubs of the tau, every one entered through its qutrit leg, and a parting at the second hub of the tau,
\begin{equation}
s(H_{2}^{\tau}) = +1 \neq w(H_{2}^{\tau}) = -1 ,
\end{equation}
because the probe attaches at the tau web at a leaf dock paired by a (2,2) link to the odd leg of that hub, so the odd-leg copy lies nearer the probe than its hub and the carry enters the hub through its qubit leg. The parting is irreducible over the three acting-site readings, a hub parting at every one. From the three lepton-seeded sinks, the electron at two sinks and the muon at one, each its own sink at wave zero and re-met by an isomorphism with its seed, by the same comparison, we obtain equality at every pendant (2,3) dock and at every hub, six of six hub readings with the multisets equal. So the sign is the side of the acting site at every hub whose odd leg is reached through its hub, and not at the one hub whose odd leg is reached ahead of its hub, which is the theorem, the open result that it is one function of the side at every hub dissolving at that re-wording. The smallest instance on the page is the electron web with its one hub, the carry from the probe dock enters the hub through its qutrit leg, the odd-leg copy reads bra, the side sign is minus one and the weight sign is minus one.

The counter-tests swapped the hub legs at the side reading, which flipped every odd-leg side sign and moved the first parting to the electron hub. read every hub-leg qubit copy as ket at the weights, which flipped the weight signs and moved the first parting to the electron hub at the reversed sign, and struck the text at its gate, which left the gate unmet. Every control read the same at both engines, among them the census of the sixty-four webs with their sites, links, hubs and ends, the orientation of the three bodies and the three sinks re-met with an isomorphism. The sides and the weights were reproduced at a second engine, equal at every hub.
\end{chproof}
\begin{figure}[h]\centering
\begin{adjustbox}{max width=\linewidth}\begin{tabular}{lllrrl}\toprule
web & hub & leg the carry enters through & side sign & weight sign & reading \\ \midrule
the electron & the hub & the qutrit leg & $-1$ & $-1$ & equal \\
the muon & the first hub & the qutrit leg & $-1$ & $-1$ & equal \\
the muon & the second hub & the qutrit leg & $-1$ & $-1$ & equal \\
the tau & the first hub & the qutrit leg & $-1$ & $-1$ & equal \\
the tau & the second hub & the qubit leg & $+1$ & $-1$ & parting \\
the tau & the third hub & the qutrit leg & $-1$ & $-1$ & equal \\ \bottomrule
\end{tabular}\end{adjustbox}
\caption{The six hubs of the lepton ladder under the carried arrow from the probe dock: the leg through which the carry enters each hub, the side sign of its odd-leg qubit copy and the weight sign of the network list there.}
\end{figure}
\par\medskip\noindent{\itshape With the sign fixed, the next chapter reads conservation as the kernel of the move matrix.}\label{narr:close:4-5}\par
\section{Conservation as the kernel of the move matrix}\label{mod:IV-6-conservation-as-the-kernel-of-the-move-matrix}
\chdesig{5-6}{4-6}{IV-6-conservation-as-the-kernel-of-the-move-matrix}{IV-6-D538128-w-apkc}
\label{narr:open:4-6}Conservation in this construction is a kernel. Theorem~74.1 gives the conserved functionals of the full move set, Theorem~74.2 the register of energy-momentum conservation as the dock, Theorem~74.3 the conserved functionals on the typed count, Theorem~74.4 the junction count at the moves, and Open result~74.5 the conserved counts as the kernel, with the kernel as the center of the Weyl algebra at the sinks, Open result~75.1, open beside them.
\begin{chtheorem}[The conserved functionals of the full move set]\label{box:IV-6-conservation-as-the-kernel-of-the-move-matrix:1}
This chapter establishes that the conserved functionals of the full move set are the left kernel of the move matrix, trivial over the grade vector, and that the counts the literature names are broken at named moves or vacuous at the sinks. The left kernel of the adopted move matrix over the grade vector has rank zero at both faces, the trivial kernel; the lepton number as the dark-triangle count is broken at a freed-at-cut column and the dart charge at the pure open-leg column.
\end{chtheorem}
\chstanding{theorem}
\chdeps{the theorems of Chapters 1-0, 1-2 to 1-5, 5-1, 5-3 to 5-4, 8-3}{Open result 75.1, Open results 76.1 to 76.3 and Theorems 78.1 to 78.3}
\begin{chconditional}[The register of energy-momentum conservation as the dock]\label{box:IV-6-conservation-as-the-kernel-of-the-move-matrix:2}
Conditional on the one-act clause and the role rule, and on the results these theorems rest on at their own standing.
The register at which energy and momentum conserve across every junction is the dock, a corollary of the dock closure; the rest-energy invariant is read at both registers and exactly one of them conserves; and the momentum transfer composes to the zero signed step along the direction register at every record, the direction functional conserved at every interaction, conditional on the one-act clause and the role rule.
\end{chconditional}
\chstanding{conditional theorem}
\chdeps{the theorems of Chapters 1-0, 1-2 to 1-5, 5-1, 5-3 to 5-4, 8-3}{Open result 75.1, Open results 76.1 to 76.3 and Theorems 78.1 to 78.3}
\begin{chconditional}[Conserved functionals on the typed count]\label{box:IV-6-conservation-as-the-kernel-of-the-move-matrix:3}
Conditional on the results these theorems rest on at their own standing, as their derivations name them.
The functionals of the typed count conserved at every census entry are the weights constant per frame dimension, the observable charges the central functionals, one $u(1)$ per mixed type; conservation is one theorem of one link, the list of twenty-four functionals placed at the corner; every acting-centre functional is conserved at every interaction of the census; and yet the move-fixed set is not the centre of the identification algebra at any link type, the known-iff-central face standing beside.
\end{chconditional}
\chstanding{conditional theorem}
\chdeps{the theorems of Chapters 1-0, 1-2 to 1-5, 5-1, 5-3 to 5-4, 8-3}{Open result 75.1, Open results 76.1 to 76.3 and Theorems 78.1 to 78.3}
\begin{chtheorem}[The junction count at the moves, restated]\label{box:IV-6-conservation-as-the-kernel-of-the-move-matrix:4}
Baryon number read as the unsigned junction count is conserved at every interaction and moved by the non-interactions alone; the structure counts at the moves are broken at the pinned release and the triangle count is broken with the join count not additive; and the lepton number as the dark-triangle count under the hub step is unequal at the freed-at-cut attachments alone.
\end{chtheorem}
\chstanding{theorem}
\chdeps{the theorems of Chapters 1-0, 1-2 to 1-5, 5-1, 5-3 to 5-4, 8-3}{Open result 75.1, Open results 76.1 to 76.3 and Theorems 78.1 to 78.3}
\begin{chwall}[The conserved counts as the kernel]\label{box:IV-6-conservation-as-the-kernel-of-the-move-matrix:5}
At the sinks the kernel over the adopted count vector has rank four at the sink edges of the initial segment and rank zero at the reduced vector; at the seeded sinks the moves are applied at every orbit representative and the kernel taken by exact arithmetic; by move class the kernels per move class nest with the full set of rank one and the electric charge parting at the duplicate; by definition line the nested kernels fall from rank nine through five to zero with the flavour pattern lost at the core family; the rank of the core-and-dock kernel moves with the placement of the hub-forming moves; the sink types at the kernel counts part at the measure and the leg count; of the ten adopted kernel functionals at the derived weights four are vacuous, two named and four unnamed; and the capacity conservation move by move parts at the first edge. The first moving cell at the placement and the kernel against the centraliser at the sinks are two further open results placed here, read at their grades.
\end{chwall}
\chstanding{open result}
\chdeps{the theorems of Chapters 1-0, 1-2 to 1-5, 5-1, 5-3 to 5-4, 8-3}{Open result 75.1, Open results 76.1 to 76.3 and Theorems 78.1 to 78.3}
The move matrix is written column by column at the grade vector and its left kernel computed by exact integer arithmetic at both faces; every admitted charge is evaluated as a vector of the grade and typed at the first column that moves it; the conserving register is read at every junction and leg against the dock register; the typed count is tabulated at every census instance and the conserved functionals solved for; the junction and triangle counts are re-met at every record network; at the sinks the deltas of every edge are keyed to the moves and the kernels computed in Hermite normal form.

\begin{chbridge}
Every continuous symmetry of an action yields a conserved charge, the theorem of Noether \cite{Noether1918}; in the Standard Model the gauge charges are Noether charges and baryon and lepton number those of accidental global symmetries \cite{Weinberg1967}; for a network of reactions with stoichiometric matrix $S$ the conserved quantities are the left kernel of $S$ \cite{Feinberg1987}. The derivation supplies the move matrix with its kernel and the dock register here, and the charges and the conservation laws are named by correspondence, the charge standing at chapters 4-4 and 4-7.
\end{chbridge}
\begin{chproof}[exact]
The premises are the moves in one table (chapter 1-4, a theorem), the content exponents of every web (chapter 1-5, a theorem) and the move matrix of the parent result at the grade vector (chapter 1-4, a theorem). The chain back to the axiom runs through chapter 1-4 and the framework chapter. From the move table, by the delta each move makes at the nine components of the grade vector (the three content exponents, the ququart count, the two orientation counts, the dark count and the two triangle counts), we obtain the move matrix at its first face, one hundred fifty-five located columns each at one integer vector,
\begin{equation} \adjustbox{max width=\dimexpr\linewidth-3em\relax}{$\displaystyle \begin{gathered} M \ \text{an integer matrix of nine rows and one hundred fifty-five columns}, \\ M_{\cdot 1} = (3,\,0,\,1,\,1,\,3,\,1,\,0,\,1,\,0)^{\mathsf T} , \end{gathered}$} \end{equation}
the first column the pure open-leg move at a qutrit, at sixteen applications at the census, and at its second face the one hundred fifty-nine columns with the four carried columns at every one of their deltas, one hundred sixty-five vectors. From the matrix, by exact integer arithmetic (the Hermite normal form of the transpose), we obtain the rank nine at both faces and the left kernel, the integer functionals annihilated by every column,
\begin{equation}
\operatorname{rank} M = 9,\qquad \ker_{\mathrm{left}} M = \{0\},\qquad \gcd(\text{maximal minors}) = 1 ,
\end{equation}
the trivial kernel at both faces, saturated, which is the theorem. From every charge named by the literature read as a vector of the grade, by its pairing with each column, we obtain its type at the first column that moves it, the lepton number as the dark-triangle count is broken at the twenty-ninth column, a cut freeing a link at a dark attachment, with the pairing minus one. The dart charge at the dock face is broken at the first column with the pairing minus six. The content exponent is broken at the first column with the pairing three, and the dock register is no functional of the grade, a register and not a count. The smallest instance on the page is the first column and the content exponent, the pure open-leg move at a qutrit raises the first content exponent by three, so the functional that reads that exponent pairs with the column at three and is not conserved,
\begin{equation}
\langle e_{a}, M_{\cdot 1}\rangle = 3 \neq 0 .
\end{equation}
The catalog of conservation laws by name, the electric charge, the lepton number, the baryon number, the color, energy and momentum, the combined parity and the parity, is mapped to a charge of the derivation where one stands and typed broken or not a functional at both faces, equal face to face at every entry.

The counter-tests struck the first column of the matrix at both faces, which moved the column counts, the two faces and the deciding cell. dropped the four carried columns, which moved the second face. reduced the grade vector to the three content exponents alone, which moved the rows, the kernels and every typing, and struck the text at its gate, which moved every cell. At the seven-component and six-component readings of the grade the kernel stays trivial. Every control read the same at both engines but one, which differed at the spelling of an empty kernel text alone, the bases equal at rank zero at every face and the control re-posed at the bases reading the same at every piece. The kernel and every charge typing were reproduced at a second engine, equal at every cell.
\end{chproof}
\begin{figure}[h]\centering
\begin{adjustbox}{max width=\linewidth}\begin{tabular}{lrrrrrrrrr}\toprule
component of the grade & $a$ & $b$ & $c$ & $q_4$ & $O_2$ & $O_3$ & $D$ & $T_H$ & $T_D$ \\ \midrule
the pure open-leg move at a qutrit & 3 & 0 & 1 & 1 & 3 & 1 & 0 & 1 & 0 \\
the pure open-leg move at a qubit & 3 & 0 & 1 & 1 & 4 & 0 & 0 & 1 & 0 \\
the cut freeing a link at a hub & 3 & 0 & 1 & 1 & 2 & 0 & 0 & 0 & 0 \\
the deletion of a freed link & 2 & 0 & 1 & 1 & 2 & 0 & 0 & 1 & 0 \\
the pairing of two hub legs & 3 & 0 & 1 & 1 & 2 & 0 & 0 & 1 & 0 \\ \bottomrule
\end{tabular}\end{adjustbox}
\caption{The move matrix at its first five forward columns: the delta of each move at the nine components of the grade vector, the three content exponents, the ququart count, the two orientation counts, the dark count and the two triangle counts; the inverse of each move is the negative column.}
\end{figure}
\section{The kernel as the centre of the Weyl algebra at the sinks}\label{mod:IV-6-D538128-w-apkc}
\begin{chwall}[The kernel as the center of the Weyl algebra at the sinks]\label{box:IV-6-D538128-w-apkc:1}
At the seven derived capped sinks, every link capped at the register of its type, the Weyl algebra of the cap records is built at every sink as the tensor product over the capped links of the pair algebras, its centre of dimension one at every sink; the commutant of the admitted group moves, the step permutation of the gate set and the holonomy permutation of every triangle, parts from that centre at every one of the seven sinks, first at the proton sink at its qutrit-qubit dock link with a commutant of dimension twelve against the centre of one, while at the reading beside, where the moves are the whole Weyl family, the commutant equals the centre everywhere; of the ten adopted kernel functionals six are vacuous at the capped set, three are Casimir elements of their factors, the dark-triangle count carrying the lepton number and two carrying the colour charge, and one is proportional to no Casimir. Noether at the cap is therefore no identity at the admitted gates, and the open result is entered with both faces named.
\end{chwall}
\chstanding{open result}
\chdeps{Definitions 1.1 to 1.6, Theorems 65.1 and 65.2 and Theorems 74.1 to 74.4 and Open result 74.5}{no theorem beyond it at this release}
The capped sinks are taken from the cap census, the Weyl algebra is built at every sink from the admitted end families with the letters as permutation and phase pairs composed at bytes and the relation of the clock and the shift checked, the center and the Casimir element of every end factor are read, the commutant of the step and holonomy gates is computed at every link and compared with the center, the filtration of the center by the link types is compared with the kernel filtration by move class, and the ten kernel functionals of the base result are evaluated at the capped sinks and matched against the Casimir elements of the eleven factor types.

\begin{chbridge}
In a gauge theory the conserved charges generate the symmetry and commute with the Hamiltonian, the Noether current of a continuous symmetry is conserved, and the operators commuting with every observable of a system span the centre of its algebra, the superselection charges; the Weyl operators of a finite register, the clock and shift matrices and their products with phases, span the full matrix algebra of the register and carry a trivial centre \cite{Noether1918,Weyl1927}. The derivation supplies the Weyl algebra of the cap records with its center and the commutant of the admitted gates here, and the conservation and the charges lepton number and color are named by correspondence.
\end{chbridge}
\begin{chproof}[finite and census]
The premises are the record phase group at every type with its Weyl relation (chapter 1-2, a theorem), the cap census at the saturated proton sinks (chapter 5-7, an open result read at its standing), the conserved functionals of the full move set (chapter 4-6, a theorem) and the Noether question at the sinks and the kernel identification (chapters 4-7 and 4-6, two open results read at their standing). The chain back to the axiom runs through chapter 1-2 and the framework chapter. The procedure. The input is the seven capped sinks of the cap census, the two proton sinks, the two stars beside, the neutron twin and the two electron controls, every one of their eighty links capped at the register of its type, and the end families at the dimensions two, three and four. At every end the letters are built as permutation and phase pairs composed at bytes, four, nine and sixteen of them, and the relation of the clock and the shift checked,
\begin{equation}
ZX = \omega\, XZ,\qquad \omega^{d} = 1 \ \text{at } d = 2, 3, 4 ,
\end{equation}
the center of every end factor read as the multiples of the identity and its Casimir element as the sum of the adjoint products of its letters beyond the identity, the scalars three, eight and fifteen, the cut register of a type the product of its two end scalars. At every sink the Weyl algebra of the cap records is the tensor product over the capped links of the pair algebras. The test applied at every capped link is whether the commutant of the admitted group moves, the step permutation of the gate set and the holonomy permutation of every triangle at its start, equals the center. The aggregation counts the sinks at which the commutant is the center and matches the ten kernel functionals of the base result against the Casimir elements of the eleven factor types. The totals at the checkpoints, the center of dimension one at seven of seven sinks. The commutant parting from the center at seven of seven sinks and equal at zero, the first parting the proton sink at its (2,3) dock link,
\begin{equation}
\dim \mathcal{A}' = 12 \quad\text{against}\quad \dim Z(\mathcal{A}) = 1 ,
\end{equation}
which is the open result. Of the ten kernel functionals six vacuous at the capped set, three Casimir elements of their factors, the dark-triangle count carrying the lepton number and two of the (3,3) pair and qutrit factors carrying the color charge, and one proportional to no Casimir, with its values minus one and zero at the seven sinks. At the reading beside, the moves read as the whole Weyl family, the commutant equals the center at every sink. The smallest instance on the page is one qubit end, its four letters are the identity, the clock, the shift and their product, the clock and the shift anticommuting, its center the multiples of the identity and its Casimir scalar three.

The counter-tests read the fifty-three seeded sinks as capped, which moved the first parting to a (2,2) link with the commutant ten. struck the qutrit factor, which moved the Casimir letter of one functional to vacuous. shifted the base result cyclically, which moved the re-meeting control. read the moves as the Weyl family, which moved the center cell to equal, and struck the text at its gate. Every control read the same at both engines, and the centers and commutants were reproduced at a second engine, equal at every link.
\chcert{80}
\end{chproof}
\par\medskip\noindent{\itshape The kernel read, the next chapter asks whether the Noether theorem of the literature has a counterpart at the sinks.}\label{narr:close:4-6}\par
\section{Noether at the sinks and the kernel identification}\label{mod:IV-7-noether-at-the-sinks-and-the-kernel-identification}
\chdesig{5-7}{4-7}{IV-7-noether-at-the-sinks-and-the-kernel-identification}{IV-7-D539098-kcfm-t1}
\label{narr:open:4-7}Open result~76.1 states the Noether theorem at the seeded sinks, Open result~76.2 the conserved charges as the center census, and Open result~76.3 the commutant spine, each an open result with its parting named.
\begin{chwall}[The Noether theorem at the seeded sinks]\label{box:IV-7-noether-at-the-sinks-and-the-kernel-identification:1}
This chapter establishes that Noether at the sinks is the kernel of the move matrix, not the commutant of the group moves and not the Cartan centraliser, which act on no count. At the fifty-three seeded sinks the kernel of the full move set restricted to every sink is computed in Hermite normal form, the commutant of the group moves is the whole count space at every sink since a group move changes no count, and the Cartan centraliser of the stabiliser construction is the whole count space as well; the three compared cell by cell agree at no sink and part at every one, first at the three-quark body.
\end{chwall}
\chstanding{open result}
\chdeps{the theorems of Chapters 1-0 to 1-2, 1-4 to 1-5, 5-1 to 5-2, 5-4, 5-6}{Chapter 5-7}
\begin{chwall}[Conserved charges as the center census]\label{box:IV-7-noether-at-the-sinks-and-the-kernel-identification:2}
The centre of the identification algebra has dimension one at every link type, and the count of independent charges equals it at three of the six types, parting first at the qubit-qubit type, while the acting centre modulo the tensor kernel equals the pair-neutral count at every type beside.
\end{chwall}
\chstanding{open result}
\chdeps{the theorems of Chapters 1-0 to 1-2, 1-4 to 1-5, 5-1 to 5-2, 5-4, 5-6}{Chapter 5-7}
\begin{chwall}[Commutant spine]\label{box:IV-7-noether-at-the-sinks-and-the-kernel-identification:3}
The commutant of the identification group at every end dimension has dimension three, spanned by the identity, the pair projector and the diagonal projector onto the matched directions, one element more than the two of the spine, and the sector reading differs at the first family and the first record.
\end{chwall}
\chstanding{open result}
\chdeps{the theorems of Chapters 1-0 to 1-2, 1-4 to 1-5, 5-1 to 5-2, 5-4, 5-6}{Chapter 5-7}
This chapter carries no theorem of its own: the kernel of the full move set restricted to every seeded sink is computed by exact integer arithmetic and compared cell by cell with the commutant of the group moves and with the Cartan centralizer carried into the count space, the center of the identification algebra is compared with the count of independent charges type by type, and the commutant of the identification group is spanned at every end dimension; every one an output of the derived results.

\begin{chbridge}
The theorem of Noether pairs every continuous symmetry of an action with a conserved charge \cite{Noether1918}; in the quantum theory a conserved charge is an operator commuting with the Hamiltonian, and the operators commuting with every element of a symmetry group form its commutant (the commutant), whose centre is spanned by the central elements of the group algebra, so that a complete set of commuting charges labels the irreducible components by the eigenvalues of a Cartan subalgebra and of the centre \cite{Hall2015,Weyl1927}. The derivation supplies the kernel of the move matrix at the sinks and the center of the identification algebra here, and the conservation named, the charges as generators and the body are named by correspondence.
\end{chbridge}
\begin{chproof}[exact and census]
The premises are the identification group at the six link types (chapter 1-2, a theorem), the conserved functionals of the full move set (chapter 4-6, a theorem), the Cartan weights at the open ends (chapter 4-4, an open result read at its standing) and the one-act clause as a theorem of the derived act (chapter 1-1, a theorem). The chain back to the axiom runs through chapter 1-2 and the framework chapter. The exact part. From the identification algebra at each of the six link types, by the dimension of its center computed at two routes, we obtain one at every type,
\begin{equation}
\dim Z(\mathcal{A}_{g}) = 1 \quad\text{at } g = (2,2),(2,3),(2,4),(3,3),(3,4),(4,4) ,
\end{equation}
and from the central functionals that are neutral on the pair, by their count at each type, we obtain the count of independent charges at the reading taken,
\begin{equation}
n_{\mathrm{pair\text{-}neutral}} = (0,\,1,\,1,\,0,\,1,\,0) \ \text{at the six types in the same order} ,
\end{equation}
one at every mixed type and zero at every homogeneous type. The two agree at three of the six types and part first at the (2,2) type, one against zero, which is the open result. Against the Cartan dimension, three, four, five, five, six and seven, the center agrees at zero of six, and the acting center, the center modulo the tensor kernel on the pair space, equals the pair-neutral count at six of six beside. The smallest instance on the page is the (2,2) type, the center of its identification algebra is the multiples of the identity, dimension one, and no central functional is neutral on the pair, so the count of independent charges is zero there. The census part. By exhaustive enumeration of the fifty-three seeded sinks (the sinks of the saturation at their seeds, the full move set restricted to every one), the property tested being whether the kernel of the restricted move set, computed in Hermite normal form, equals the commutant of the group moves and the Cartan centralizer carried into the count space. The counts by kind, the commutant of the group moves the whole count space at every sink, since a group move changes no count, the Cartan centralizer the whole count space as well, and the three agreeing cell by cell at no sink and parting at every one, the first parting element the three-quark body. So the conserved functionals at the sinks are the kernel and neither of the other two, which act on no count. At every end dimension the commutant of the identification group has dimension three, spanned by the identity, the pair projector and the diagonal projector onto the matched directions, one element more than the two of the spine.

The counter-tests of the derivation stood at their pinned letters and the controls held. The center dimensions and the charge counts were reproduced at a second route, equal at every type, and the kernels at the sinks were compared with the commutants at both faces.
\chcert{53}
\end{chproof}
\begin{figure}[h]\centering
\begin{adjustbox}{max width=\linewidth}\begin{tabular}{lrrrr}\toprule
link type & ends & center of the identification algebra & pair-neutral charges & Cartan dimension \\ \midrule
(2,2) & 2, 2 & 1 & 0 & 3 \\
(2,3) & 2, 3 & 1 & 1 & 4 \\
(2,4) & 2, 4 & 1 & 1 & 5 \\
(3,3) & 3, 3 & 1 & 0 & 5 \\
(3,4) & 3, 4 & 1 & 1 & 6 \\
(4,4) & 4, 4 & 1 & 0 & 7 \\ \bottomrule
\end{tabular}\end{adjustbox}
\caption{The identification algebra at the six link types: the dimension of its center, the count of central functionals neutral on the pair and the Cartan dimension beside; the center and the count part first at the (2,2) type.}
\end{figure}
\section{The kernel as the commutant of the full move set}\label{mod:IV-7-D539098-kcfm-t1}
A set of operators acting irreducibly on a finite-dimensional space has a commutant consisting of the scalars alone (the lemma of Schur), so the conserved quantities of a dynamics generated by such a set are the Casimir elements of its algebra, central and constant on every irreducible block; the centre of the algebra of the moves is the algebra of the conserved charges \cite{Schur1905,Weyl1931,Noether1918}. The kernel identification is introduced at section~\ref{mod:IV-7-noether-at-the-sinks-and-the-kernel-identification} and conservation as the kernel of the move matrix at section~\ref{mod:IV-6-conservation-as-the-kernel-of-the-move-matrix}.
At every capped link of the seven capped sinks, the group generated by the one-end moves at either end together with the step transport and the holonomy of every triangle at its start has a commutant of dimension one at each of the six genres, so the commutant of the full move set equals the centre at 7 of 7 capped sinks and the conserved charges are the centre alone.
The seven sinks carry 80 capped links at 80 of 80. The generated groups at the ends are those of the universality result, of orders 576, 10368 and 6144 with the Hadamard read over two among them; adjoining the step transport of the gate-set result and the triangle holonomies, the commutant on the link tensor algebra has dimension 1 at every genre, against 10, 12, 24, 45, 32 and 136 for the step and holonomy gates alone, the Weyl-algebra reading being reproduced at every link. The generated group at every end acts irreducibly on its end register at 3 of 3, every operator commuting with every move is a scalar at every capped sink, the three Casimir functionals of the naming table, the dark-triangle factor at 728, the (3,3) link factor at 80 and the qutrit factor at 8, each the letter count less one, are elements of the centre at 3 of 3, and the standing obstruction, that the commutant of the step and holonomy gates alone parts from the centre, is superseded at the full move set.
The result is a theorem, conditional on the results it cites at their own standing: the kernel is the centre at every capped sink, the first face read and the second found at no sink; no comparison with observation is made here.
What stays open is a case in which the one-end moves at some end generate a group with a commutant beyond the scalars on its end register, found at none; the Weyl-algebra obstruction (section~\ref{mod:IV-6-D538128-w-apkc}) stands at its own question, superseded here at the full move set.
\par\medskip\noindent{\itshape From conservation the thread passes to the identity it implies at one link, and to the anomaly.}\label{narr:close:4-7}\par
\section{The Ward identity and the anomaly sums}\label{mod:IV-8-ward-identity-and-the-anomaly-sums}
\chdesig{5-8}{4-8}{IV-8-ward-identity-and-the-anomaly-sums}{IV-8-D540512-aedd-t1}
\label{narr:open:4-8}Theorem~78.1 reads the Ward identity from conservation as one link, Theorem~78.2 the reading of the records washed, and Theorem~78.3 the anomaly read at the body part, with Theorem~79.1 placing the anomaly with every body at its own sink.
\begin{chconditional}[The Ward identity from conservation as one link]\label{box:IV-8-ward-identity-and-the-anomaly-sums:1}
Conditional on the reading of the records it names, a record read as the one link subdivided at the dimension of each of its read ends in sequence at the first admitted ordering, and on the results it rests on at their own standing.
This chapter derives the Ward identity from conservation as one link, washes its reading of the records, and defines the anomaly read at the body part with its cross term positive. At the typed-count reading the vertex dressing of a configuration differs from its count by the sum over its link pairs of the typed count the two ends of each link would carry if open, so the vertex equals the count exactly when the per-dimension sums are constant, the one conserved line of the one-link theorem; at the subdivision of a dressed link, one read and three joins, the vertex through the new link pairs equals the line through the undivided link at every functional, the count unchanged. The result is conditional on the reading of the records it names.
\end{chconditional}
\chstanding{conditional theorem}
\chdeps{Definitions 1.1 to 1.6, Theorems 9.1 to 9.7 and Open result 9.8, Theorems 18.1 to 18.4 and Open results 18.5 and 18.6, Theorems 28.1 and 28.2, Open results 28.3 and 28.4 and Theorem 28.5, Open results 66.1 to 66.3, Theorem 66.4 and Open result 66.5, Open results 72.1 to 72.3, Theorem 72.4 and Open result 72.5, Theorems 74.1 to 74.4 and Open result 74.5 and Open results 241.1 and 241.2, Theorems 241.3 and 241.4 and Open results 241.5 to 241.7}{Theorem 70.1, Theorem 79.1, Theorem 80.1 and Open result 95.1}
\begin{chconditional}[The reading of the records washed]\label{box:IV-8-ward-identity-and-the-anomaly-sums:2}
Conditional on the results this theorem rests on at their own standing, as its derivation names them.
That reading is washed: the identity holds at every one of the forty-one readings of the records the definitions admit, and the path of four one-act moves is forced at the pair state, so the adopted reading is derived.
\end{chconditional}
\chstanding{conditional theorem}
\chdeps{Definitions 1.1 to 1.6, Theorems 9.1 to 9.7 and Open result 9.8, Theorems 18.1 to 18.4 and Open results 18.5 and 18.6, Theorems 28.1 and 28.2, Open results 28.3 and 28.4 and Theorem 28.5, Open results 66.1 to 66.3, Theorem 66.4 and Open result 66.5, Open results 72.1 to 72.3, Theorem 72.4 and Open result 72.5, Theorems 74.1 to 74.4 and Open result 74.5 and Open results 241.1 and 241.2, Theorems 241.3 and 241.4 and Open results 241.5 to 241.7}{Theorem 70.1, Theorem 79.1, Theorem 80.1 and Open result 95.1}
\begin{chtheorem}[The anomaly read at the body part]\label{box:IV-8-ward-identity-and-the-anomaly-sums:3}
With the adopted placement applied at the measurement web of the anomalous-dipole result, every record of the direct sum over records is classified into the body, the cross records across the join and the apparatus self-terms, and the body-part anomaly is defined with the orders separated: at order one the unit itself, the dock-struck first power of the electron, sixty-seven, with the sign plus; at order two an exact element of the field with the sign plus and a nonzero cross term; the order-three term zero. The hub spelling, the muon and the tau are printed beside; every number is an exact object of the derivation, compared with nothing here. The anomaly as the ratio of the two reads at the electron is an open result placed at this chapter and read at its grade.
\end{chtheorem}
\chstanding{theorem}
\chdeps{Definitions 1.1 to 1.6, Theorems 9.1 to 9.7 and Open result 9.8, Theorems 18.1 to 18.4 and Open results 18.5 and 18.6, Theorems 28.1 and 28.2, Open results 28.3 and 28.4 and Theorem 28.5, Open results 66.1 to 66.3, Theorem 66.4 and Open result 66.5, Open results 72.1 to 72.3, Theorem 72.4 and Open result 72.5, Theorems 74.1 to 74.4 and Open result 74.5 and Open results 241.1 and 241.2, Theorems 241.3 and 241.4 and Open results 241.5 to 241.7}{Theorem 70.1, Theorem 79.1, Theorem 80.1 and Open result 95.1}
The vertex dressing of every configuration of the one hundred seventy-two records is written at the typed-count reading as the count plus the sum of the open-end counts over its link pairs, the vanishing of that sum is shown equivalent to the constancy of the per-dimension sums, and the subdivision path is telescoped at every one-link, census and listed path; every reading of the records the definitions admit is enumerated and the identity checked at every reading and record; and every record of the direct sum over records of the measurement web is classified by the placement theorem with the sums taken at every order.

\begin{chbridge}
The Ward identity of quantum electrodynamics relates the vertex function to the propagator, $k_\mu \Gamma^\mu(p+k, p) = S^{-1}(p+k) - S^{-1}(p)$, so that the vertex and the wave-function renormalisations are equal, $Z_1 = Z_2$, a consequence of current conservation \cite{Ward1950,Takahashi1957}; the anomalous dipole factor of the electron is $a_e = (g-2)/2$ with the one-loop term $\alpha/2\pi$ of Schwinger \cite{Schwinger1949}; and the axial anomaly of a chiral current cancels in the Standard Model at every generation when the linear, cubic and mixed sums of the hypercharges over the sixteen Weyl components vanish \cite{Adler1969,BellJackiw1969,Bouchiat1972}. The derivation supplies the vertex dressing and the count of a configuration with their identity at every path here, and the Ward identity, the anomaly, and the dipole factor are named by correspondence.
\end{chbridge}
\begin{chproof}[census and exact]
The premises are the one-act clause as a theorem of the derived act (chapter 1-1, a theorem), conservation as one theorem of one link (chapter 4-6, a theorem), the explicit two-point function with its one hundred seventy-two records (chapter 5-1, a theorem) and the dressing at the listed face with its two thousand six hundred thirty-three dressed links (chapter 1-3, a theorem). The chain back to the axiom runs through chapter 4-6 and the framework chapter. The exact part. From a configuration of the one link at the typed-count reading, by the action of a generator on its state factor by factor, the open copies at their typed weight, every link pair at the typed count its two ends would carry if open and every joining pair zero, we obtain the vertex dressing as the count plus a sum over the link pairs,
\begin{equation}
V(\text{configuration}) = C + \sum_{\text{link pairs}} \big(w(\text{end}_1) + w(\text{end}_2)\big) ,
\end{equation}
so the vertex equals the count exactly when the per-dimension sums are constant, the one conserved line of the one-link theorem. At the one link the vertex, the count and the line read zero, zero and zero at the polarity plus, while at the open result generator the vertex reads minus twenty-four against the count minus twelve, the identity failing there. From the subdivision of a dressed link, one read and three joins, by the telescoping of the dressing through the new link pairs, we obtain the vertex through the subdivided link equal to the line through the undivided link at every functional of the frozen list, the count unchanged. The census part. By exhaustive enumeration of the paths of the subdivision (the set generated by every admitted ordering at every inserted dimension and both polarities), the identity holds at forty-eight of forty-eight paths at the one link, at one thousand six hundred fifty-three of one thousand six hundred fifty-three census subdivisions over thirteen thousand two hundred twenty-four paths, at two thousand six hundred thirty-three of two thousand six hundred thirty-three dressed links of the listed face with the count zero at every one, and at their subdivisions
\begin{equation}
63192 \text{ of } 63192 \text{ paths},\qquad \text{the open result generator at } 0 \text{ of } 63192 ,
\end{equation}
and at the one hundred seventy-two records of the one link at one hundred seventy-two of one hundred seventy-two, the first failing path none, the central generator commuting with every letter at every record. So the identity is a theorem at the reading named, which is the theorem. The smallest instance on the page is the one link at the polarity plus, its vertex, its count and its line all read zero, and its one subdivision path telescopes to the same zero.

The counter-tests of the derivation stood at their pinned letters. Two controls mis-posed at the spelling of the list, the link types read with the larger end first at twelve lines and the copy names sorted as text, were re-spelled by the review at two routes with every cell unmoved, disclosed. The identity was reproduced at a second engine, equal at every path.
\chcert{63192}
\end{chproof}
\section{The anomaly with every body at its own sink}\label{mod:IV-8-D538315-aebs-t1}
\begin{chtheorem}[The anomaly with every body at its own sink]\label{box:IV-8-D538315-aebs-t1:1}
Every body of the anomaly grid replaced by its own sink carries a self-term at its sink: the electron and the muon sinks are the identity and the tau is read at the tau-form sinks, and the ratio of the two reads re-cut at every derived pairing of the base grid is unmoved at the faces of the orders one, two and three and moved at the whole face alone, by the tau-form sink alone, the moved cells named per body at every reading of that sink.
\end{chtheorem}
\chstanding{theorem}
\chdeps{Definitions 1.1 to 1.6, Theorems 78.1 to 78.3 and Theorem 141.1}{no theorem beyond it at this release}
The sinks of the three leptons are taken from the saturation, the self-term is read at every sink, the ratio of the two reads is re-cut at every derived pairing of the base grid at the four faces, and the cells that move are named per body; the counter-tests are met and the controls agree at both engines, one control re-spelled by the review at the cells of the item where an empty value was spelled at the word of each engine.

\begin{chbridge}
The anomalous magnetic dipole factor of the electron is the correction to the Dirac value of its gyromagnetic ratio, computed in quantum electrodynamics as a series in the fine-structure constant whose first term is the Schwinger term, with the contributions of the heavier leptons and of the hadrons entering at higher orders through the vacuum polarization; the ratio of the muon and electron anomalies tests the universality of the lepton couplings \cite{Schwinger1948,Kinoshita1990}. The derivation supplies the self-term of a body at its sink and the ratio of the two reads here, and the anomaly and the dipole factor are named by correspondence.
\end{chbridge}
\begin{chproof}[finite and census]
The premises are the ratio of the two reads at the electron (chapter 4-8, a theorem), the saturation with the sinks of the three leptons (chapter 1-5, a theorem) and the lepton-form sinks of the tau-sink saturation (an open result placed at no chapter of this compendium, read at its standing). The chain back to the axiom runs through chapter 4-8 and the framework chapter. The procedure. The input is the base grid of the anomaly, its one hundred forty-four derived pairings at the four faces, the two spellings and the three parts, and the sinks of the three leptons, the electron and the muon at their own sinks, which are the identity, and the tau at the tau-form sinks read at the seed face with the electron and star faces beside. At every sink the self-term of the body is read at the order two, and at every pairing the ratio of the two reads is re-cut with every body replaced by its own sink. The test applied at each pairing and face is whether the ratio moves under the replacement, and the aggregation counts the pairings moved and names the moved cells per body. The totals at the checkpoints,
\begin{equation} \adjustbox{max width=\dimexpr\linewidth-3em\relax}{$\displaystyle \begin{gathered} \text{pairings moved} = 24 \text{ of } 144 \ \text{at every reading}, \\ \text{moved at the faces of the orders one, two and three} = 0 , \end{gathered}$} \end{equation}
every moved pairing at the whole face and the moved body the tau alone, its self-term at the order two at the sink against the listed copy, the entry order five against five at the seed and star faces and four against five at the electron face, the first moved cell the hub spelling at the body part over the one-link office. The electron and the muon replacements the identity, and no body whose sink carries no self-term, the join defined at the valence bound at every replaced body. So every body replaced by its own sink carries a self-term there and the ratio is unmoved at the three lower faces and moved at the whole face by the tau alone, with the moved cells named, which is the theorem. The smallest instance on the page is the electron, its sink is its own listed web, the replacement is the identity, and the ratio of its two reads is unmoved at every pairing.

The counter-tests struck the self-term at the seed face, which moved one hundred fifty-five cells beginning at the sink cell. struck the dock spelling of the body part, which moved seventy-five. struck the derived pairings, which moved two hundred fifty-seven, and struck the text at its gate, which moved every cell. Every control read the same at both engines but one, re-spelled by the review at the cells of the item where an empty value was spelled at the word of each engine, and the ratios were reproduced at a second engine, equal at every pairing.
\chcert{144}
\end{chproof}
\section{The images moved by a recorded cell}\label{mod:IV-8-D540512-aedd-t1}
\begin{chtheorem}[The images moved by a recorded cell]\label{box:IV-8-D540512-aedd-t1:1}
The one hundred forty-four plus forty-eight images of the magnetic-read row, re-read at the derived definition of the open end and the far-end dressing, are moved by a recorded cell at every image with the ratio exact: the images are listed at the frame of one hundred ninety-two at every variant, the one hundred forty-four at the base faces and thirty-six at the orders four and five over the one-link, two-dock and whole-sink offices, one hundred eighty of one hundred eighty exact elements of the field at every variant, the twelve at the symmetric part beyond the base faces carrying the word of the row and no image; under the recorded definition the smaller end of every unequal link is open and both ends of an equal link are washed, under the derived definition the larger end is open, the two letters equal at the twenty ends of the electron web at the recorded sets and the slot of the read engine the larger end at one hundred nine of one hundred sixteen unequal links; the electron carries one hub with its in-leg at (3,4) and three legs at (2,4), every far end open under the recorded definition with the trace one and washed under the derived with the traces three, two, two and two, the derived dressing eight at the lepton legs against the recorded dressing three; at the derived polarity the order one of the probe moves from $\tfrac{16195}{1934917632} + \tfrac{695}{161243136}\sqrt{2} + \tfrac{601}{161243136}\sqrt{3} + \tfrac{751}{322486272}\sqrt{6}$ to $\tfrac{21385}{1934917632} + \tfrac{1225}{161243136}\sqrt{2} + \tfrac{1}{157464}\sqrt{3} + \tfrac{1057}{241864704}\sqrt{6}$, the first power sixty-seven unmoved and the closure of the two-dock office moved to $\tfrac{2904369707}{95551488000} + \tfrac{18271}{911250}\sqrt{2} + \tfrac{10254403}{39813120000}\sqrt{3}$, the office $\tfrac{428933534001707}{95551488000} + \tfrac{18271}{911250}\sqrt{2} + \tfrac{10254403}{39813120000}\sqrt{3}$; the images are moved at one hundred eighty of one hundred eighty at the seed variant and alike at the three beside, inexact none, unmoved none, the moving cells recorded. The first face reads yes and the second face, unmoved at every image or the first parting image named, reads no at every variant.
\end{chtheorem}
\chstanding{theorem}
\chdeps{Definitions 1.1 to 1.6, Theorems 78.1 to 78.3, Theorems 185.1 to 185.12 and Theorem 195.1}{no theorem beyond it at this release}
The images of the magnetic-read row are re-derived from the captures of the self-terms engines over the mass cells by locator; the letter at every end is read under the recorded and the derived definitions with the polarity census; the joined web of every body is rebuilt at the docking face and the self-terms re-read at both polarities; the two-dock office is re-read at the derived polarity, and every image is compared with its cell at the base row.

\begin{chbridge}
The anomalous magnetic coupling of the electron is a series in the coupling whose coefficients beyond the leading term depend on the convention of the subtraction, so that a change of convention moves every higher image while the leading term stands \cite{Schwinger1948,Petermann1957}. The derivation supplies the images at both definitions and the ratio at every one here, and the coupling is named by correspondence.
\end{chbridge}
\begin{chproof}[finite and exact]
The premises are the magnetic-read row at its images (chapter 4-8, a theorem), the larger-end open result at the derived definition (chapter 1-7, an open result), the read-cap premises row at the open sets (chapter 1-7, an open result), the dressing row at the hub legs (chapter 8-1, a theorem), the self-terms row at the captures (chapter 4-8, a theorem) and the body-alone office of the electron (chapter 8-2, a theorem). The chain back to the axiom runs through chapter 4-8 and the framework chapter. The procedure. The input is the one hundred ninety-two images of the base row at four variants, the thirteen bodies, the photon word and the tau-form sink webs with their ends, and the captures of the self-terms engines. The enumeration rule re-derives every exact image from the self-terms at every body summed by the depth rule over the mass cells by the field inverse, reads the letter at every end under both definitions, rebuilds the joined web of every body at the docking face as the self-terms engines build it, re-reads the self-terms and the two-dock office at the derived polarity, and compares every image with its cell at the base row; the test at each image is whether it moves with the ratio exact. From the self-terms at the recorded polarity, equal to the captures at fifty-two of fifty-two bodies and the sums at every order, we obtain the images re-met at one hundred eighty of one hundred eighty at every variant, the order one of the probe at the dock spelling over the one-link office
\begin{equation}
\tfrac{16195}{8685845250048} + \tfrac{695}{723820437504}\sqrt{2} + \tfrac{601}{723820437504}\sqrt{3} + \tfrac{751}{1447640875008}\sqrt{6} ,
\end{equation}
and from the derived polarity the same image moved to $\tfrac{21385}{8685845250048} + \tfrac{1225}{723820437504}\sqrt{2} + \tfrac{1}{706855896}\sqrt{3} + \tfrac{1057}{1085730656256}\sqrt{6}$, the tau-form sink at the order two moved from $\tfrac{7017711173}{1074954240000} + \tfrac{2521}{583200}\sqrt{2} + \tfrac{1252321}{358318080000}\sqrt{3}$ to $\tfrac{3542099269}{537477120000} + \tfrac{128087}{29160000}\sqrt{2} + \tfrac{2723}{746496000}\sqrt{3}$ entering at the order five, the sums moved at every order one to five and at the whole. The counts by kind: the letters at the twenty ends of the electron web equal to the recorded open set and the derived set with the read-open letter at the far end at twenty of twenty, the three leptons re-met against the dressing row at twelve of twelve; the polarity census one hundred sixteen unequal links with the slot the larger end at one hundred nine, the seven others the carrier words and the photon word at pendant qutrit ends of degree one and the tau-form sink at two legs; the traces of the electron legs three, two, two and two under the derived definition, re-met against the dressing row at twenty-seven of twenty-seven; the two-dock office at the recorded polarity equal to the body-alone row at ten of ten with the one-link office four thousand four hundred eighty-nine; the closed form of every passage at both polarities equal to the read engine at two hundred twenty-eight of two hundred twenty-eight; the images moved at one hundred eighty of one hundred eighty at every variant, inexact none. So every image is moved by a recorded cell with the ratio exact, which is the theorem. The smallest instance on the page is the first image, the order one of the probe at the dock spelling, the two cells displayed above.

The counter-tests struck the derived reading back to the recorded reading, which moved eleven cells with the face among them and every image unmoved. A second struck the far-end dressing back to the recorded reading at the derived reading, which moved the symmetric-part image at the body alone. A third struck the closure sector at every record space, which moved sixteen cells with the electron self-term among them, and a fourth struck the gate, which left every cell absent. Every control read the same at both engines, fifteen of fifteen, and the images were reproduced at a second engine of a different construction, byte-identical at every run.
\chcert{228}
\end{chproof}
\begin{figure}[h]\centering
\begin{adjustbox}{max width=\linewidth}\begin{tabular}{lll}\toprule
image & at the recorded polarity & at the derived polarity \\ \midrule
the order one of the probe at the dock spelling & $\tfrac{16195}{8685845250048} + \tfrac{695}{723820437504}\sqrt{2} + \cdots$ & $\tfrac{21385}{8685845250048} + \tfrac{1225}{723820437504}\sqrt{2} + \cdots$ \\
the first power of the electron & $67$ & $67$ \\
the closure of the two-dock office & the recorded cell & $\tfrac{2904369707}{95551488000} + \tfrac{18271}{911250}\sqrt{2} + \tfrac{10254403}{39813120000}\sqrt{3}$ \\
the dressing of the electron at the lepton legs & $3$ & $8$ \\ \bottomrule
\end{tabular}\end{adjustbox}
\caption{Four cells of the re-read at the two polarities, the first power unmoved and every image moved by the recorded closure and dressing cells.}
\end{figure}
\par\medskip\noindent{\itshape The next chapter states the one rule every read obeys and the table of words it produces.}\label{narr:close:4-8}\par
\section{The instrument tables of this Part}\label{mod:V-L-the-instrument-tables-at-the-repository}
\chdesig{5-L}{}{V-L-the-instrument-tables-at-the-repository}{V-L-the-instrument-tables-at-the-repository}
\begin{chtable}[The instrument tables of this Part]\label{box:V-L-the-instrument-tables-at-the-repository:1}
This Part lists 12 instrument chapters, printed among the instruments at the end of this compendium as chapters~\ref{mod:IV-8b-D540345-the-lagrangian-table-at-the-twelve-slots} to~\ref{mod:IV-14-bell-bound-column}, and the table beneath names each one with its key cells.
\end{chtable}
\chstanding{instrument}
\chdeps{Definitions 1.1 to 1.6}{no theorem beyond it at this release}
The listing moves no cell and states no result of its own. Each instrument stands at the repository as its chapter stood at the previous release of this compendium, and its tables are read there.
\begin{chbridge}
The tables listed here are instruments of the derivation, and no name of the literature is attached at this listing.
\end{chbridge}
\begin{chproof}[instrument]
The listing in full. For each instrument chapter, printed among the instruments at the end of this compendium, the table gives its title, the result boxes it carries, the standing of each box with its citation, and the count of objects its certificate sentence names where it prints one.
{\small
\begin{longtable}{p{38mm}p{63mm}p{30mm}p{14mm}}
\hline
chapter & results & standing & objects \\
\hline
\endhead
The twelve slots: the table & The Lagrangian table at the twelve slots & instrument & -- \\
The twelve slots: the table & The twelve slots at both definitions & instrument & -- \\
The three registers and the six pairings & The six genres as the pairs of the three registers & instrument & -- \\
The three registers and the six pairings & Exactly six pairs of the three registers & theorem & 1691 \\
The three registers and the six pairings & Three registers reached and nothing else & theorem & 36628 \\
The body membership and the quantum numbers & The body membership census at the six cells & instrument & -- \\
The body membership and the quantum numbers & The anatomy atlas of the recorded sinks & instrument & -- \\
The SM table: the 75 entries with their status & The standard-model entries at their grades; theorem & instrument; theorem & -- \\
The SM table: the 75 entries with their status & the text of the chapter & see its text & -- \\
The SM table: the 75 entries with their status & The grades of the textbook corpus & instrument & -- \\
The face groups of the six cells & The face groups of the six cells & instrument & -- \\
The bell bound column & The Bell bound column & instrument & -- \\
\hline
\end{longtable}
}
The rule that produced it. Every chapter of this Part whose first standing beyond a placeholder reads instrument is printed among the instruments at the end of this compendium with its tables, and its key cells are the titles of its boxes, its standing lines and its certificate count, read from its text as it stood at the previous release.
No derivation stands in this chapter; the tables stand at the repository.
\emph{Check.} Every file listed was re-read at the repository and read the same as its recorded copy, and the listing was re-cut from the files by a second route.
Nothing stands open at this listing. At this release three results of this Part are re-graded at their chapters: the quantum numbers of every body (the chapter 5-11) re-posed at the re-cut catalogue, the unconserved sites (the chapter 5-15) re-posed at the first parting site, the frame at every end (the chapter 5-16) re-posed at its own link.
\end{chproof}
\label{chend:V-L-the-instrument-tables-at-the-repository}
\section{The ququart as the register of the first power}\label{mod:IV-8f-D542554-w-msrg}
\chdesig{5-8f}{4-8f}{IV-8f-D542554-w-msrg}{IV-8f-D542664-nctg-t1}
\begin{chwall}[The ququart as the register of the first power]\label{box:IV-8f-D542554-w-msrg:1}
The ququart is not the register of the first power over the whole recorded census under the body-alone rule as the parents apply it, the first body parting the photon word: the office at every body and word of the census at the adopted reading is an exact element of the field of root two and root three at fifty-four of fifty-four reads with every site dimension one of the three registers, the Higgs word at the first power sixteen and the office $\tfrac{46656001681}{182250000}$, the W word at twenty-eight and $\tfrac{2857690241}{3645000} + \tfrac{98}{50625}\sqrt{2}$, the electron at two hundred one and $\tfrac{2\ 144\ 666\ 856\ 770\ 839}{53084160000} + \tfrac{661}{4050}\sqrt{2} + \tfrac{5977541}{2654208000}\sqrt{3}$, the tau at three thousand nine hundred sixty-nine, the tau sink at four thousand one hundred eighty-five, the cross-core web at thirty-six and the top word at two hundred twenty-eight; the fifty reads touching the ququart are nonzero at fifty of fifty, every one with a ququart hub or a ququart end at a (2,4), (3,4) or (4,4) link; the four reads within the qubit and the qutrit are at the office zero at two of four, the one-link web of the body-alone row and the up docked body with every link struck, and nonzero at the other two, the photon word, six sites in a chain of five (2,3) links with the two pendant docks struck and the three inner links unstruck at the act six each, at the first power eighteen and the office three hundred twenty-four, and the gluon word, four qutrits in a chain of three (3,3) links with nothing struck, at the first power twenty-seven and the office $\tfrac{20639133397}{28311552} + \tfrac{1675}{7077888}\sqrt{3}$. The first face, the office zero exactly within the qubit and the qutrit and nonzero exactly at the ququart, reads no at the adopted reading and at none of the two hundred forty readings of the grid, and the second face reads yes with the first body parting the photon word at three hundred twenty-four against zero, the partings two of fifty-four, the two neutral carrier words and no body among them.
\end{chwall}
\chstanding{open result}
\chdeps{Definitions 1.1 to 1.6, Theorems 185.1 to 185.12, Theorem 195.1 and Theorems 200.1 to 200.3 and Open result 200.4}{no theorem beyond it at this release}
The census of bodies and words is transcribed at fifty-four reads with every parent read re-met at its cell by locator; the office of every read is computed under the body-alone rule at the adopted reading and at a grid of two hundred forty readings; the reads are sorted by whether they touch the ququart, and the first face is tested at every reading with the first parting named in the census order.

\begin{chbridge}
In the Standard Model the particles coupled to the Higgs doublet acquire a mass proportional to that coupling, while the photon and the gluons, carrying no such coupling, remain massless, so that the Higgs sector is the register of mass \cite{Weinberg1967,Higgs1964}. The derivation supplies the sort of the fifty-four reads by their offices here, and the register is named by correspondence.
\end{chbridge}
\begin{chproof}[finite and census and exact]
The premises are the body-alone office rule (chapter 8-2, a theorem), the Higgs and carrier offices (chapter 6-10, a theorem), the cross-core offices (chapter 7-9, a theorem), the tau web row (chapter 1-5, a theorem), the body membership census (chapter 4-8e, an instrument) and the mass unit with the six dressings (chapter 8-1, a theorem). The chain back to the axiom runs through chapter 8-2 and the framework chapter. By exhaustive enumeration of the fifty-four reads of the recorded census, thirty-six distinct structures with the electron two-dock web entering at six sources, the muon and the tau at four each, the top word at two and the Higgs, W and Z words at two each, the property tested at each read being whether its office under the body-alone rule is zero and whether it touches the ququart. From the photon word, the two pendant docks struck and the three inner (2,3) links unstruck at the act six each, by the first power and the closure, we obtain
\begin{equation}
p = 6 + 6 + 6 = 18 , \qquad c = 0 , \qquad \omega = p^{2} + c = 324 ,
\end{equation}
no dark link at the adopted dark set so the closure zero and the office the first power squared, nonzero; and from the gluon word, three (3,3) links at the act nine each with nothing struck, the first power twenty-seven, the pair part seven hundred twenty-nine and the closure $\tfrac{11989}{28311552} + \tfrac{1675}{7077888}\sqrt{3}$ over three dark links, the office $\tfrac{20639133397}{28311552} + \tfrac{1675}{7077888}\sqrt{3}$, nonzero. The counts by kind: the reads exact at fifty-four of fifty-four with every parent read equal to its recorded cell by locator; the reads touching the ququart nonzero at fifty of fifty, the Higgs word and the leg piece, the W and Z words, the closed-genre objects, the free genres, the electron, the muon, the tau, the proton, the neutron, the deuteron, the two stars, the three cross-core webs, the three lepton sinks, the top word and the four docked bodies; the reads within the qubit and the qutrit zero at two of four; the first face at none of the two hundred forty readings of the grid, the photon and gluon chains nonzero at every act reading since their inner links are unstruck at every one. So the sort holds at every body of the census and at the Higgs, W and Z words and parts at the two neutral carrier words alone, which is the open result. The smallest instance on the page is the one-link web of the body-alone row: one (2,3) link with its dock struck at the pendant site, the first power zero and the office zero.

The counter-tests struck the first-named ququart site of the Higgs word with its one link, which moved its office to zero. A second inserted a ququart site at the one-link web by a (3,4) link, which moved its office from zero to nonzero and its read into the ququart sort. A third read the office at the first power alone, which moved one hundred two cells with the office of the Higgs word among them, a fourth hid the (3,4) cells, which moved the electron, and a fifth struck the gate, which left every cell absent. Every control read the same at both engines, fifteen of fifteen, and the grid was reproduced at a second engine of a different construction, byte-identical at every run.
\chcert{54}
\end{chproof}
\begin{figure}[h]\centering
\begin{adjustbox}{max width=\linewidth}\begin{tabular}{llrl}\toprule
the read within the qubit and the qutrit & its links & first power & the office \\ \midrule
the one-link web of the body-alone row & one (2,3) link, the dock struck & 0 & $0$ \\
the up docked body & two (2,3) docks, both struck & 0 & $0$ \\
the photon word & five (2,3) links, three unstruck & 18 & $324$ \\
the gluon word & three (3,3) links, nothing struck & 27 & $\tfrac{20639133397}{28311552} + \tfrac{1675}{7077888}\sqrt{3}$ \\ \bottomrule
\end{tabular}\end{adjustbox}
\caption{The four reads of the census within the qubit and the qutrit at the adopted reading, the office zero at two and nonzero at the two neutral carrier words, the photon word the first parting.}
\end{figure}
\section{Three, a theorem of the valence}\label{mod:IV-8f-D542664-nctg-t1}
\begin{chtheorem}[Three, a theorem of the valence]\label{box:IV-8f-D542664-nctg-t1:1}
The generation count three is a theorem of the valence law at the three core sites of the closed triangle: a rung is a (3,4) in-leg at a hub, a site of dimension four at degree three or more, with its far end the qutrit, and the core sites are the qutrits on the closed (3,3) links; at every recorded body the rungs are the electron one at its hub, the muon two, the tau three with every core site at room zero, the proton three at the pair-less hub from the junction qutrits, the neutron four and the tau sink beside three, the exponents of the passage zero, two, three, two, three and three with the dressings one, nine, twenty-seven, nine, twenty-seven and twenty-seven equal to the unpassed-in-legs cells at thirty-six of thirty-six and the in-legs, hubs and counts equal to the legs cells at eighteen of eighteen; the saturation at the (3,4) cell at the adopted reading, the docks struck, the attach sites the core sites, the unit the hub assembly with its pair and odd leg, the site bound twenty-four and the wave cap eight, terminates at the rung three at every lepton seed, the bare core at the wave three over eight webs, the electron at the wave two over four, the muon at the wave one over two and the tau at the wave zero over one web with no admitting site, every one exact and at the same rung at the widened bound forty-eight, the proton and the neutron stopping at six at this reading and at three and four with the docks kept, printed beside; the count of rungs at termination is one count, three, at every lepton seed, the chain from the electron at the attach sites isomorphic to the recorded muon and tau webs at three of three, and the rung order from the bare core equal to the attach-site-basis order. The first face reads yes with the recorded lepton rungs one, two and three within the count, the greatest equal to it, and a further rung admitted at none of the core sites of the tau; the second face, another count or no bound, reads no at the adopted reading with the first parting none, and at the readings beside the (2,4) cell gives the count zero with the tau at the rung nine, the struck rung gives differing counts and the lifted bound gives no bound, every one printed beside.
\end{chtheorem}
\chstanding{theorem}
\chdeps{Definitions 1.1 to 1.6, Theorems 185.1 to 185.12 and Open results 245.1 and 245.2, Theorem 245.3, Open results 245.4 to 245.8 and Theorem 245.9}{no theorem beyond it at this release}
The rungs of the (3,4) cell are read at every recorded body and re-met at the unpassed-in-legs, legs and generation-ladder cells; the saturation at the (3,4) cell is run from the bare core and the three lepton seeds at the adopted reading and at the readings beside with every new hub assembly checked at its degree and valence; the count of rungs at termination is read, the chain from the electron is matched to the muon and tau webs by backtracking, and the parent rows are read at their faces by locator.

\begin{chbridge}
The Standard Model has three generations of quarks and leptons, identical in their gauge quantum numbers and differing in their couplings to the Higgs, and three is the least count at which the quark mixing carries a phase \cite{Kobayashi1973,Weinberg1995}. The derivation supplies the count of rungs at the termination here, and the three generations are named by correspondence.
\end{chbridge}
\begin{chproof}[finite and exact]
The premises are the unpassed in-legs at the five bodies (chapter 8-1, a theorem), the generation ladder and the generation bound (chapter 4-8f, a theorem and an open result), the attach-site basis at the core sites (chapter 4-8f, an open result), the termination of every seeded saturation at its counting bound (chapter 1-5, a theorem), the tau sink row at its labels (chapter 1-5, an open result) and the valence law at every site (chapter 1-3, a theorem). The chain back to the axiom runs through chapter 1-3 and the framework chapter. By exhaustive enumeration of the hub assemblies the saturation at the (3,4) cell admits from the bare core and the three lepton seeds, the bare core at eight webs over the attach keys, the electron at four, the muon at two and the tau at one, thirty-three new hub assemblies every one at degree equal to its dimension with one in-leg and twenty-seven reached webs within the valence bound at every site, the property tested at each web being whether a further rung is admitted at a core site with room. From the rungs at the recorded bodies, by the count of (3,4) in-legs at hubs, we obtain
\begin{equation}
r_{e} = 1 , \quad r_{\mu} = 2 , \quad r_{\tau} = 3 , \quad r_{p} = 3 , \quad r_{n} = 4 , \qquad k = 0, 2, 3, 2, 3 ,\ \ \text{the dressings } 3^{k} = 1, 9, 27, 9, 27 ,
\end{equation}
equal to the unpassed-in-legs cells at thirty-six of thirty-six and to the ladder, the total and the exponent of the generation rows at the three leptons. From the saturation at the adopted reading, every step at a site with room alone, we obtain the termination at the rung three at every lepton seed, the bare core stopped at the wave three over the eight webs of the attach keys, the electron at the wave two over four, the muon at the wave one over two and the tau at the wave zero over one web, the count
\begin{equation}
N_{\text{rungs}} = 3 \quad \text{at the core, the electron, the muon and the tau} ,
\end{equation}
exact and unmoved at the widened bound forty-eight, the hadron seeds at six with the docks displaced beside. The counts by kind: the new hub assemblies at degree equal to dimension with one in-leg thirty-three of thirty-three and the reached webs within the valence twenty-seven of twenty-seven; the chain from the electron at the attach sites in order isomorphic to the recorded muon and tau webs at three of three by backtracking at both engines, the rung order equal to the attach-site-basis order; the counting bound of the termination row recomputed as an integer equal to its cells at the site bounds four, five, six and twelve; the fourth label of the tau-sink row at this reading recounted equal with the nine sinks beyond the third rung typed seven at four or more attach sites and two at a repeated attach site and none at the closed-triangle form; the parent rows at their faces by locator thirteen of thirteen. So the count of rungs at termination is three, the recorded lepton rungs within it with the greatest equal to it, which is the theorem. The smallest instance on the page is the tau seed: three rungs at its three hubs, every core site at room zero, stopped at the wave zero with no admitting site.

The counter-tests widened the budget, which left the count three, the termination rungs and the deciding cell unmoved. A second read the (2,4) cell in place of the (3,4), which moved the rungs to three, six and nine at the leptons with the count zero and the second face. A third struck a rung, which moved the bare core, the electron and the muon seeds to two with the tau at three, a fourth struck the termination bound with the site bound lifted, which left every seed unstopped at the wave cap with no bound, and a fifth struck the gate, which left every cell struck. Every control read the same at both engines, twenty of twenty, and the counts were reproduced at a second engine of a different construction, byte-identical at every run.
\chcert{36}
\end{chproof}
\begin{figure}[h]\centering
\begin{adjustbox}{max width=\linewidth}\begin{tabular}{lrrrl}\toprule
the seed & rungs at the seed & exponent & dressing & the termination at the (3,4) cell \\ \midrule
the bare core & 0 & & & the rung 3 at the wave 3 over 8 webs \\
the electron & 1 & 0 & 1 & the rung 3 at the wave 2 over 4 webs \\
the muon & 2 & 2 & 9 & the rung 3 at the wave 1 over 2 webs \\
the tau & 3 & 3 & 27 & the rung 3 at the wave 0 over 1 web \\
the proton & 3 & 2 & 9 & 6 at this reading, 3 with the docks kept, beside \\
the neutron & 4 & 3 & 27 & 6 at this reading, 4 with the docks kept, beside \\ \bottomrule
\end{tabular}\end{adjustbox}
\caption{The rungs of the (3,4) cell at every recorded body with the exponent and dressing of its passage and the termination of the saturation at the adopted reading, one count, three, at every lepton seed.}
\end{figure}
\section{The one rule of a read and the word table}\label{mod:IV-9-one-rule-of-a-read-and-the-word-table}
\chdesig{5-9}{4-9}{IV-9-one-rule-of-a-read-and-the-word-table}{IV-9-one-rule-of-a-read-and-the-word-table}
\label{narr:open:4-9}Open result~84.1 states the one rule of a read and the word table, Open result~84.2 the passage web of the cross read against the W word, and Open result~84.3 the Z read as the two-line read of the W word, every one an open result that Part VIII returns to.
\begin{chwall}[The one rule of a read and the word table]\label{box:IV-9-one-rule-of-a-read-and-the-word-table:1}
This chapter states the one opening rule of the act at a sink as four derived clauses, the strike rule with the dock rule, the leg count, the hub-leg rule and the orientation law, and finds the sixteen words reached by the union of the conventions and by no one rule: applied at the sixty seeded sinks with every end lettered, the rule reaches nine of the sixteen words at the one face under the adopted convention, nine under the end-letter convention and nine with the orientation law, thirteen at the union of the three, and the sixteen at the wide union alone; the first word reached by the union alone is the docked up row, reached by the docked down body whose two docks hang at one qutrit site, where the clause the one rule lacks is the multiplicity per qutrit site, a clause of no derived result.
\end{chwall}
\chstanding{open result}
\chdeps{the theorems of Chapters 1-0, 1-2 to 1-5, 5-1, 5-3 to 5-5}{no theorem beyond it at this release}
\begin{chwall}[The passage web of the cross read against the W word]\label{box:IV-9-one-rule-of-a-read-and-the-word-table:2}
This open result is stated once, as Open result~159.2 of the chapter The W and the Z and their masses, and this chapter cites it there without restating it.
\end{chwall}
\chstanding{open result}
\chdeps{the theorems of Chapters 1-0, 1-2 to 1-5, 5-1, 5-3 to 5-5}{no theorem beyond it at this release}
\begin{chwall}[The Z read as the two-line read of the W word]\label{box:IV-9-one-rule-of-a-read-and-the-word-table:3}
This open result is stated once, as Open result~159.1 of the chapter The W and the Z and their masses, and this chapter cites it there without restating it.
\end{chwall}
\chstanding{open result}
\chdeps{the theorems of Chapters 1-0, 1-2 to 1-5, 5-1, 5-3 to 5-5}{no theorem beyond it at this release}
This chapter carries no theorem of its own: the four clauses of the opening rule are applied at every seeded sink with every end lettered and the words reached are counted per convention and per face; the passage web of the cross read is built from the derived cross and compared with the W word of the network list at the link types, the rank and the sites; and the two-line read of the W word is computed at its six forms with the two routes of the carriers result compared at every one.

\begin{chbridge}
In the electroweak theory the neutral carriers are the mixtures $Z = \cos\theta_W W^3 - \sin\theta_W B$ and $A = \sin\theta_W W^3 + \cos\theta_W B$ of the third weak field and the hypercharge field, the W and Z masses standing at the ratio $m_W/m_Z = \cos\theta_W$ at tree level, and the charged current couples the two members of every doublet \cite{Glashow1961,Weinberg1967}; the content of a generation is one table of sixteen Weyl components, every one a word in the weights of the gauge group. The derivation supplies the opening rule of the act with the words it reaches at the sinks and the two-line read of the W word here, and the W and the Z, the weak angle, and the mass are named by correspondence, the weak angle standing at the fence alone.
\end{chbridge}
\begin{chproof}[finite]
The premises are the strike rule with the dock rule, the leg count, the hub-leg rule and the orientation law as four derived clauses (chapters 1-3, 1-4 and 4-5, read at their standing), the saturation with its seeded sinks (chapter 1-5, a theorem) and the sixteen words of the network list as a comparison target and no input (chapter 4-1, a theorem). The chain back to the axiom runs through chapter 1-4 and the framework chapter. The procedure. The input is the sixty seeded sinks with every end lettered, fifty-three of fifty-three, three of three and four of four at their three kinds, their nine hundred seventy-eight colored ends with none open under the strike clause, and the one opening rule of the act at a sink stated as four clauses, no clause a selection. The rule is applied at every sink at its one face and the words reached are read under three conventions, the label convention, the end-letter convention and the end-letter convention through the orientation law. The test applied at every sink and convention is which of the sixteen words of the list are reached. The aggregation counts the words reached per convention, per face and at the unions. The totals at the checkpoints,
\begin{equation} \adjustbox{max width=\dimexpr\linewidth-3em\relax}{$\displaystyle \begin{gathered} \text{words reached at the one face: } 9,\ 9,\ 9 \text{ of } 16 \text{ under the three conventions}, \\ \text{their union } 13 \text{ of } 16, \qquad \text{the wide union } 16 \text{ of } 16 , \end{gathered}$} \end{equation}
the leg count of the second clause equal to the derived column at sixty of sixty, and the words reached per sink at the one face forty-nine sinks at one word, two at none, two at two, three at three and four at six. No face under one convention reaches the sixteen at zero of sixty-three face-convention pairs, the most nine of sixteen. The first word of the frozen order unreached at the one face under every single convention is the first right row of the docked up body, reached at the wide union at the six per-site faces alone by the docked down body, whose two docks hang at the one qutrit site at the far end of its hub leg, the hub-leg clause per dock link opens that qutrit end twice, the anti-triplet of weight two at the fundamental sign and no row, while the multiplicity per qutrit site opens it once, the triplet of weight four and the docked up right rows. So the clause the one rule lacks is the multiplicity per qutrit site, a clause of no derived result, which is the open result. The smallest instance on the page is that docked down body, two (2,3) docks at one qutrit site, the end opened twice by the rule and once by the multiplicity, two words against one.

The counter-tests struck the strike clause, the leg-count clause, the hub-leg clause and the orientation clause in turn, which moved one thousand three hundred eleven, one hundred sixty, six hundred sixty-five and two hundred sixty-eight cells of the three thousand and some each, and struck the text at its gate, which left the gate unmet. Every control read the same at both engines, and the words reached were reproduced at a second engine, equal at every sink.
\chcert{978}
\end{chproof}
\par\medskip\noindent{\itshape The Part closes on the table that records the standing of every entry of the model.}\label{narr:close:4-9}\par
\section{The quantum numbers of every body as the end counts of its membership}\label{mod:IV-11-quantum-numbers-of-every-body-as-the-end-counts-of-its-membership}
\chdesig{5-11}{4-11}{IV-11-quantum-numbers-of-every-body-as-the-end-counts-of-its-membership}{IV-11-quantum-numbers-of-every-body-as-the-end-counts-of-its-membership}
\begin{chwall}[The quantum numbers of every body as the end counts of its membership]\label{box:IV-11-quantum-numbers-of-every-body-as-the-end-counts-of-its-membership:1}
The end counts of the membership give the recorded quantum numbers at every resolved body and part from the recorded cells at a named first body: the open ends of every body at the pendant reading, a site of degree one at the link spelling, are counted by register and, at the qubit ends, by the cell they sit at, with the (2,2) links apart, at seventy-eight of eighty-three bodies, the link spelling re-derived from the recorded webs and the four parent tables with the dimensions propagated to a fixed point and equal to the census at seventy-eight of eighty-three, the five uncounted bodies carrying one unequal link or two of one genre at a shared site whose spelling fixes no register at either end; at the electron web the open qubit ends are three, two at (2,3) cells, the docks, and one at a (2,4) cell, the kept leaf, every letter bra, with no open qutrit or ququart end and one (2,2) link, at the up quark docked web two open qubit ends at (2,3) cells of opposite letters, at the W word no open qubit end and two open ququart ends, at the photon word one open qubit end and one open qutrit end. The quantities at every resolved body are printed: the doublets and the twice-spin set, the statistics by the parity, odd at thirty-nine and even at thirty-nine of seventy-eight, the colour content, one at seventy-three, three at four and three by two at one of seventy-eight at the pendant reading, the triality, the weak doublet count, zero at twenty-three, one at twenty-two and two at seventeen, the hypercharge at two readings, the Dirac form, yes at sixteen of seventy-eight, the letters of the open ends with the self-conjugacy and the face by name against the expected face, met at forty-five of sixty-seven; at the electron the doublets three, the statistics odd, the colour one, the weak doublet count one with the su(2) content two, the hypercharge four and nine at the two readings, the Dirac form yes and the face Dirac by name equal to the expected. The open results at their named partings are re-met at twenty-two of twenty-eight, nine as instances, thirteen located only and six differing with both cells named, the first the representation census row at the electron, the su(2) content two by the one open (2,4) qubit end against the cell one at the dock count. The first face, every recorded quantum number at every recorded body, reads no at seventy-eight of eighty-three counted, sixty-seven of eighty-three printed and twenty-two of twenty-eight re-met; the second face reads yes with the first parting, in the order of the table by membership, the (3,4) one-link free genre of the Higgs and carrier row at the end counts themselves, undefined by the adopted reading because its spelling fixes no register at either end, and among the printed bodies the electron at the weak content, two against one, and the doublet content, three against one, the up quark docked web at the statistics and the face, and the photon word at the statistics, every one recorded beside.
Re-posed at the re-cut catalogue of webs: the end counts part from the recorded cells at the same first body, the (3,4) one-link free genre of the Higgs and carrier row, a body of a parent table alone whose spelling fixes no register at either end, undefined at the re-cut catalogue because the re-cut touches no parent table; the forty-five kept bodies at the forty-nine kept webs re-derived there equal the recorded cells at every compared key, forty-five of forty-five, the fifteen struck bodies at the uniform-chain rows uncounted, the first at the open qubit ends against the recorded three, the twenty-three bodies of the parent tables carried, the bodies counted sixty-three of eighty-three; the twenty-eight partings re-posed there re-meet at twenty-two of twenty-eight, every one standing at the re-cut catalogue and at its certificates, the first differing the su(2) content at the electron, two against one.
\end{chwall}
\chstanding{open result}
\chdeps{Definitions 1.1 to 1.6, Theorem 57.1 and Open result 57.2, Open results 66.1 to 66.3, Theorem 66.4 and Open result 66.5 and Theorem 124.1}{no theorem beyond it at this release}
The link spelling of every body is re-derived from the recorded webs and the parent tables with the dimensions propagated to a fixed point and compared with the census; the open ends are counted at the pendant reading by register and cell; the quantities of the item are read at every resolved body at every reading and printed; the twenty-eight open results are re-met at their named partings as instances, as located or as differing with both cells, and the first parting is named in the order of the table by membership.

\begin{chbridge}
Every particle of the Standard Model is labeled by its spin, its statistics, its color, its weak isospin and its hypercharge, and the spin-statistics theorem ties the first two, so that a table of the fields is a table of these labels \cite{Pauli1940,Weinberg1995}. The derivation supplies the end counts at every resolved body and the partings from the recorded cells here, and the labels are named by correspondence.
\end{chbridge}
\begin{chproof}[finite and census]
The premises are the body membership census (chapter 4-8e, an instrument), the representation census at its rows (chapter 4-3, theorems), the spin-label row at the open qubit ends (chapter 3-5, an open result), the colour rows at the hub-leg reading (chapter 4-4, open results), the Dirac-form open results at the docked link (chapter 3-5, open results), the face-at-the-ports row (chapter 7-9, an open result) and the stabiliser row at the central scalars (chapter 5-1, an open result). The chain back to the axiom runs through chapter 1-2 and the framework chapter. By exhaustive enumeration of the eighty-three bodies of the census with their link spellings over the sixty-four recorded webs and the four parent tables, the property tested at each body being whether its spelling fixes the register of every end, whether its end counts give every quantity of the item and whether each open result naming it is re-met at its parting. From the spelling of the electron web, by the propagation of the dimensions to a fixed point, we obtain the open ends
\begin{equation}
n_{2} = 3 \ \text{(two at (2,3), one at (2,4), every letter bra)} , \qquad n_{3} = 0 , \qquad n_{4} = 0 , \qquad \text{the (2,2) links } 1 ,
\end{equation}
the sorted spelling equal to the census at that body, and from the counts by the rules of the item the doublets three, the twice-spin set three and one, the statistics odd, the colour one at the pendant reading and three by three at the odd reading with the triality zero at both, the weak doublet count one with the su(2) content two, the hypercharge four at the set reading and nine at the links reading, the Dirac form yes and the face Dirac by name, the typing of the two docks spatial and the leaf internal. The counts by kind: the bodies counted seventy-eight of eighty-three, the five uncounted the (3,4) free genre and leg piece, the (2,4) free genre, the port-only cross-core pair and the one-link web of the body-alone row; the bodies printed sixty-seven of eighty-three, the first unprinted the re-keyed cross-core web with the conjugation letters undefined at its open ends; the statistics odd at thirty-nine and even at thirty-nine; the colour at the hub-leg reading equal to the colour rows at fifty-three of fifty-three webs; the spin-label table equal at the open qubit ends at sixty-four of sixty-four; the open results re-met at twenty-two of twenty-eight, the instances at nine, the colour rows at the two-dock down body, the spin-label row at the up quark web, the photon word and the W word, the up-type-quark row at the typing of the dock, the face-at-the-ports row and the lepton-number row at the ports of the three cross-core webs, the located-only at thirteen, every one a parting outside the end counts, and the differing at six, the representation census at the electron, two against one, the odd-cycle row at the W word, zero against two, and at the electron, three against one, the doublet-swap row at the gluon word, zero against two, and the doublet-kind row at the electron. So the end counts give the quantities at every resolved body and re-meet the colour rows, the spin-label ends and the port orientations as instances while parting from the recorded cells at the doublet content of the leptons and carriers, at the weak content of the electron and at the registers of the one-link bodies, the first parting the (3,4) free genre, which is the open result. The smallest instance on the page is the one-link web of the body-alone row: one (2,3) link with its dock struck at the pendant site, its spelling fixing the register of the pendant end alone.

The counter-tests struck the kept leaf of the electron web, which moved seventeen cells. A second added a (2,4) end at the up quark web, which moved twenty-five cells. A third applied the swap at the (2,3) cell of the electron web alone, which moved two cells, a fourth took a membership row with no open end, the Higgs word closed, which moved four cells, and a fifth struck the gate, which left the gate unmet. The controls read the same at both engines at eleven of fifteen, the four differing at the five unresolved bodies, at nine of ten offices, at the engine parity of the typed set and at the cross-diff, the first two mis-posed and the third one parting between the engines at two of four thousand five hundred four shared keys, every one disclosed.
\chcert{83}
\end{chproof}
\section{The diagonal self couplings in one form}\label{mod:IV-13-diagonal-self-couplings-in-one-form}
\chdesig{5-13}{4-13}{IV-13-diagonal-self-couplings-in-one-form}{IV-13-diagonal-self-couplings-in-one-form}
\begin{chtheorem}[The diagonal self couplings in one form]\label{box:IV-13-diagonal-self-couplings-in-one-form:1}
The three diagonal self-couplings, the self term of the (2,2) genre, the bracket of the (3,3) genre at the gluon pair ket and the quartic of the (4,4) genre at the vacuum register, are three cells of one law at three genres, the count law read at the diagonal, and no second law enters: at (2,2) the count law gives the traceless letters three and three, the sectors twenty-four, eighteen, twelve and three at the gains one, two, two and four, fifty-seven records, the Born count ninety-six, the registers sixteen and nine, the second-order read $\tfrac{5}{216} + \tfrac{5}{216}\sqrt{2}$, the first power four and the pair part sixteen, and the increment of the count metric along the recorded chain is ninety-six, ninety-six, ninety-six, equal to the Born count at three of three with the second difference zero, while the register tower multiplies itself by the whole of the (2,3) genre per step; at (3,3) the count law gives the letters eight and eight, the sectors one hundred eighty-four, one hundred twenty-eight, seventy-two and twenty-eight at the gains one, three, three and nine, four hundred twelve records, the Born count one thousand thirty-six, the registers eighty-one and sixty-four, the read $\tfrac{67}{82944} + \tfrac{25}{41472}\sqrt{3}$, the first power nine and the pair part eighty-one, and the identification algebra of dimension seventeen closes at two hundred eighty-nine of two hundred eighty-nine brackets with the far-end commutator pairs twenty-eight equal to the closure sector and twenty-seven nonzero, the stabiliser of the pair ket of dimension nine; at (4,4) the count law gives the letters fifteen and fifteen, the sectors six hundred sixty, four hundred fifty, two hundred forty and one hundred five at the gains one, four, four and sixteen, one thousand four hundred fifty-five records, the Born count five thousand one hundred, the registers two hundred fifty-six and two hundred twenty-five, the read $\tfrac{41}{216000}$, the first power sixteen, the pair part two hundred fifty-six and the one-link office $\tfrac{46656001681}{182250000}$, and the quartic, the office over the vacuum register six thousand five hundred sixty-one squared, is $\tfrac{46656001681}{7\ 845\ 264\ 902\ 250\ 000}$, equal to the Higgs-quartic and Higgs-and-carrier rows. The one column at the three genres is the Born count ninety-six, one thousand thirty-six and five thousand one hundred, the closure sector three, twenty-eight and one hundred five, the second-order read at each, the whole register sixteen, eighty-one and two hundred fifty-six, and the stabiliser of dimension four, nine and sixteen with its centre one line not acting on the pair ket; the re-meets against the recorded rows are equal at sixty-five of sixty-five. The first face reads yes at the three genres and the second face, the first parting, reads no, the first parting none.
\end{chtheorem}
\chstanding{theorem}
\chdeps{Definitions 1.1 to 1.6, Theorems 90.1 to 90.4, Theorem 124.1, Theorems 185.1 to 185.12 and Theorems 200.1 to 200.3 and Open result 200.4}{no theorem beyond it at this release}
The count law is evaluated at the three diagonal genres from the traceless letters and the sectors; the (2,2) self term is re-met at the count metric along the recorded chain and at the increment operator of the curvature open result; the (3,3) bracket is built in the recorded basis of the stabiliser row and closed with its commutator pairs counted against the closure sector; the (4,4) quartic is formed as the office over the vacuum register squared and re-met at the quartic rows; the one column is printed at the three genres and the parent rows are read at their faces.

\begin{chbridge}
The self-couplings of the Standard Model, the self-energy of the photon field, the cubic and quartic self-interactions of the gluons fixed by the structure constants, and the quartic of the Higgs potential, are each set by one coupling of their sector and by no second parameter \cite{YangMills1954,Higgs1964}. The derivation supplies the one column at the three diagonal genres here, and the self-couplings are named by correspondence.
\end{chbridge}
\begin{chproof}[census and exact]
The premises are the count law at the six genres (chapter 8-1, a theorem), the two metrics at the depth (chapter 10-0, an open result), the curvature as the Laplacian washed over the ladder (chapter 10-1, an open result), the stabiliser of the pair ket (chapter 5-1, an open result), the gluon pair ket at its record (chapter 4-4, a theorem), the Higgs quartic and the Higgs and carrier offices (chapter 6-10, theorems) and the six genres at their read orders (chapter 4-8d, an instrument). The chain back to the axiom runs through chapter 8-1 and the framework chapter. By exhaustive enumeration of the brackets of the identification algebras at the three diagonal genres, forty-nine at (2,2), two hundred eighty-nine at (3,3) and nine hundred sixty-one at (4,4), with the Jacobi identity at thirty-five, six hundred eighty and four thousand four hundred ninety-five triples, every one in scope, the property tested at each bracket being closure in the recorded basis and at each triple the identity. From the count law at the diagonal genre of end dimension $d$, by the sectors at their gains, we obtain the Born count and the second-order read,
\begin{equation} \adjustbox{max width=\dimexpr\linewidth-3em\relax}{$\displaystyle \begin{gathered} B_{2} = 96 , \quad B_{3} = 1036 , \quad B_{4} = 5100 , \\ G_{2}(2,2) = \tfrac{5}{216} + \tfrac{5}{216}\sqrt{2} , \quad G_{2}(3,3) = \tfrac{67}{82944} + \tfrac{25}{41472}\sqrt{3} , \quad G_{2}(4,4) = \tfrac{41}{216000} , \end{gathered}$} \end{equation}
the closure sectors three, twenty-eight and one hundred five at the gain $d^{2}$. From the recorded chain of the two-metrics open result, by the difference, the count metric ninety-six, one hundred ninety-two, two hundred eighty-eight and three hundred eighty-four at the counts one to four has the increments ninety-six, ninety-six and ninety-six, equal to the Born count, with the second difference zero and the dependence on the count none, while the register tower twenty-four multiplied by thirty-six, then thirty-one thousand one hundred four, one million one hundred nineteen thousand seven hundred forty-four and forty million three hundred ten thousand seven hundred eighty-four has its increments linear in the tower, the tower over the count metric nine, one hundred sixty-two, three thousand eight hundred eighty-eight and one hundred four thousand nine hundred seventy-six equal to the two-metrics cells at four of four; the second-difference increment operator at an inserted (2,2) link on its three-site support, two one-link Laplacians added and one removed, has the trace two, the rank two, the spectrum minus one, zero and three and the row sums zero, equal to the curvature open result increment at twenty-four of twenty-four. From the (3,3) identification algebra in the recorded basis, by the brackets, we obtain the closure at two hundred eighty-nine of two hundred eighty-nine with one hundred eight nonzero at the ordered pairs, every structure constant among minus two, minus one, one and two, the far-end commutator pairs twenty-eight equal to the closure sector with twenty-seven nonzero, the two diagonal letters alone commuting, the stabiliser of dimension nine and the centre $1;1;1;-1;-1;-1$ at the pair scalar zero. From the (4,4) office at one link, by the recorded vacuum register, the (3,3) whole register eighty-one to the order exponent two,
\begin{equation}
\lambda = \frac{\omega_{(4,4)}}{V_{1}^{2}} = \frac{46656001681}{182250000 \cdot 6561^{2}} = \frac{46656001681}{7\ 845\ 264\ 902\ 250\ 000} ,
\end{equation}
equal to the Higgs-quartic hub-face quartic and the Higgs-and-carrier quartic with the inverse identity, and at the register beside eighty-one the quartic $\tfrac{46656001681}{1195742250000}$ equal to its cell. The counts by kind: the re-meets at (2,2) sixteen of sixteen, at (3,3) thirty-one of thirty-one and at (4,4) eighteen of eighteen, sixty-five of sixty-five; the parent rows at their deciding cells consistent at nine of nine. So the (2,2) term sits at the Born count, the (3,3) term at the closure sector and the order-two read, and the (4,4) term at the order-two read squared at the pair part over the vacuum register, one law at three genres, which is the theorem. The smallest instance on the page is the (2,2) genre: the closure sector three, the three commutator pairs of the traceless letters at the gain four.

The counter-tests put the (2,3) pair in place of the (3,3), which moved forty-two cells with the sectors seventy-nine, forty-eight, forty-two and three and the centre acting at the pair scalar one. A second struck the bracket, which moved fifteen cells with the far-end nonzero pairs zero. A third read the quartic at the shifted register, which moved thirteen cells, a fourth struck the gate, which left every cell struck, and a fifth added the (5,5) pair compared with nothing, which moved eight cells with the three deciding cells unmoved. Every control read the same at both engines, sixteen of sixteen, and the column was reproduced at a second engine of a different construction, byte-identical at every run.
\chcert{4495}
\end{chproof}
\begin{figure}[h]\centering
\begin{adjustbox}{max width=\linewidth}\begin{tabular}{lrrr}\toprule
the cell of the one column & (2,2) & (3,3) & (4,4) \\ \midrule
the Born count & 96 & 1036 & 5100 \\
the closure sector & 3 & 28 & 105 \\
the whole register & 16 & 81 & 256 \\
the first power & 4 & 9 & 16 \\
the stabiliser of the pair ket & 4 & 9 & 16 \\
the self-coupling it carries & the self term & the bracket & the quartic \\ \bottomrule
\end{tabular}\end{adjustbox}
\caption{The one column of the count law at the three diagonal genres, the self term at the Born count, the bracket at the closure sector and the quartic at the order-two read squared over the vacuum register.}
\end{figure}
\section{The unconserved sites as the hub acts}\label{mod:IV-15-unconserved-sites-as-the-hub-acts}
\chdesig{5-15}{4-15}{IV-15-unconserved-sites-as-the-hub-acts}{IV-15-unconserved-sites-as-the-hub-acts}
\begin{chwall}[The unconserved sites as the hub acts]\label{box:IV-15-unconserved-sites-as-the-hub-acts:1}
The unconserved sites of the continuity equation are not the sites at which the hub acts of the depth rule act, the first site parting at the universe web named: the recorded webs are the twenty-one webs of the ladder of the curvature row, three readings at the rungs zero to six rebuilt from the closure rows, the universe web at four hundred forty-four sites and five hundred fifty-one links; the office density at every site is the sum over the links at the site of the product of the end dimensions, eight thousand eight hundred thirty-six over the universe web, re-met at the linear row at forty-two of forty-two keys; the chains are the components of the (2,2) links, eighty-five of one link at the universe web and three at the closure-site body; a chain is conserved when the density is one value along it, and at the universe web seventy-eight of eighty-five chains are conserved and seven unconserved, every one a bare qubit at the density four against the third leg of a hub unit at the density twelve, the fourteen unconserved sites the item figures fourteen and seven re-met, with every chain conserved at every web of the first two readings and the same fourteen sites at the seven webs of the third; the hub acts of the depth rule are located at every web by the patterns of the move graph, the dock a site of dimension two at degree one whose neighbour is of dimension three and the hub form a hub of dimension four at degree three with one neighbour of dimension three and two paired legs of dimension two at degree two, and at the universe web the hub-act sites are one, the one dock, with no hub form since the seventy-five sites of the hub dimension stand every one at degree four; the two site sets compared at every one of the twenty-one webs at five readings part at the universe web with none in common at the adopted reading, the first parting site the third leg of a hub unit, unconserved and not a hub act, and are empty and equal at the closure-site body at the rungs one to six and at the open-end chain at the odd rungs. At the seven item webs, the universe web and the webs at the rungs one and two of every reading, the first face, the unconserved sites equal to the hub-act sites at every web, reads no with three of seven equal where both sets are empty, and the second face reads yes with the first parting at the universe web at that third leg, at the density twelve against its bare qubit at four; every one of the fourteen unconserved sites is a third leg of a hub unit or the partner of a third leg, fourteen of fourteen within the one hundred one sites of that reading with the eighty-seven beside conserved, so the sources of the continuity equation sit at the third legs of the hub step and their partners and not at the hub acts themselves.
Re-posed at the re-cut catalogue of webs as the parent material: the unconserved sites and the hub-act sites of the hub-forming move of the depth rule are equal at three of the seven webs of the question, where both sets are empty, and part at the first, the universe web, at its first unconserved site, a qubit at degree two on a hub and its bare partner at the density twelve against four along their chain, the third leg of a hub unit and not a hub act, a hub act at the readings that count the unit, the third leg or the third-leg pair and not at the readings that count the hub form or the hub step, the parting above re-met at its site sets at every key; the re-cut catalogue beside at forty-nine webs is equal at twenty-one of forty-nine, its first parting at the dock of the docked up quark, a hub act and no unconserved site, its unconserved sites at eleven of its webs every one a third leg of a hub unit or a leg of a representation web.
\end{chwall}
\chstanding{open result}
\chdeps{Definitions 1.1 to 1.6, Theorems 90.1 to 90.4, Open result 92.1 and Theorem 92.2, Theorems 105.1 and 105.2, Open results 105.3 to 105.5 and Theorem 105.6 and Open results 241.1 and 241.2, Theorems 241.3 and 241.4 and Open results 241.5 to 241.7}{no theorem beyond it at this release}
The ladder is rebuilt from the closure rows and re-met at the nonlinear row; the office density is read at every site and the chains of (2,2) links are traced with their conservation; the hub acts of the depth rule are located at every web by the patterns of the move graph at four readings beside; the two site sets are compared at every web and reading with the first parting site named in hand-off order, and the nonlinear row is transcribed by key and re-met at every web.

\begin{chbridge}
A continuity equation has sources where the current is not conserved, and in a gauge theory the sources of the field sit at the charges, not at the vertices where the interaction acts \cite{Maxwell1865,Noether1918}. The derivation supplies the two site sets at every web and their first parting here, and the sources are named by correspondence.
\end{chbridge}
\begin{chproof}[finite and exact]
The premises are the curvature as the Laplacian washed over the ladder (chapter 10-1, an open result), the nonlinear row at its conservation counts (chapter 10-2, an open result), the linear row at its density (chapter 10-2, an open result), the depth rule with its move table (chapter 10-4, an open result) and the move-graph open result at its edges (chapter 1-4, an open result). The chain back to the axiom runs through chapter 10-1 and the framework chapter. By exhaustive enumeration of the twenty-one webs of the ladder at their sites, four hundred forty-four at the universe web, and of the five readings of the hub acts at every web, the property tested at each web being whether the set of unconserved sites equals the set of hub-act sites and, where not, which site parts first in hand-off order. From the density at every site, by the one value along each chain, we obtain at the universe web the seven unconserved chains, each a bare qubit at the density four joined to the third leg of a hub unit at the density twelve,
\begin{equation}
\rho(\text{bare qubit}) = 4 , \qquad \rho(\text{third leg}) = 12 , \qquad \text{the unconserved sites } 14 = 2 \cdot 7 ,
\end{equation}
re-met at the nonlinear row at every web at two hundred seventy-three of two hundred seventy-three keys and at the operator reading, the sites where the row-sum-zero operator applied to the density is not zero, at twenty-one of twenty-one webs. From the move table of twenty-eight moves closed under the inverse rule at fourteen of fourteen pairs, the hub-forming move adding two sites and the join keeping them, by the patterns at the pinned dimensions, we obtain the hub-act sites at the universe web, one dock and no hub form, and at the readings beside the anchors of the seventy-five hub units, the two hundred ninety-two sites the hub steps add, the seventy-three third legs and the one hundred one third legs with their partners, the patterns re-met at the fourteen roster webs of the move graph at forty-two of forty-two cells and the move-graph open result re-met at its sixty-eight thousand forty-four edges by kind. The counts by kind: the two sets at the universe web at the adopted reading parting with none in common and the first parting site unconserved and not a hub act; at the anchors parting with none in common; at the unit sites seven in common; at the third legs seven in common and seven unconserved alone; at the third legs with their partners fourteen in common, none unconserved alone and eighty-seven hub-act alone; at the closure-site body at the rung zero and the open-end chain at the even rungs parting at the one dock with no unconserved site; at the other webs both sets empty; the item webs equal at three of seven; all twenty-one webs equal at nine of twenty-one. So the sources sit at the third legs of the hub step and their partners and not at the hub acts, which is the open result. The smallest instance on the page is one unconserved chain of the universe web: a bare qubit at the density four joined to a third leg at the density twelve, the two densities not one value.

The counter-tests struck the hub-forming pattern at the dock dimension, which moved the hub-act count at the universe web from one to zero. A second took the (3,4) chains in place of the (2,2), which moved the unconserved sites from fourteen to one hundred thirty-one. A third struck the density uniform, which moved the unconserved chains to zero, a fourth struck the gate, which left every cell absent, and a fifth switched the hub-act reading to the unit sites at the deciding cell, which moved the first parting site to a unit site. Every control read the same at both engines, twelve of twelve, and the sets were reproduced at a second engine of a different construction, byte-identical at every run.
\chcert{21}
\end{chproof}
\begin{figure}[h]\centering
\begin{adjustbox}{max width=\linewidth}\begin{tabular}{lrrrl}\toprule
the reading of the hub acts at the universe web & in common & unconserved alone & hub-act alone & the first parting site \\ \midrule
the hub-forming move, the adopted reading & 0 & 14 & 1 & a third leg, unconserved and not a hub act \\
the hub step at its anchors & 0 & 14 & 75 & an anchor, conserved and a hub act \\
the four sites every hub step adds & 7 & 7 & 285 & a unit site, conserved and a hub act \\
the third legs & 7 & 7 & 66 & a third leg, conserved and a hub act \\
the third legs with their partners & 14 & 0 & 87 & a third leg, conserved and a hub act \\ \bottomrule
\end{tabular}\end{adjustbox}
\caption{The unconserved sites of the universe web against the hub-act sites at the five readings, parting at every one, with the fourteen unconserved sites all among the third legs and their partners.}
\end{figure}
\section{The frame at every end as the stabiliser of its dimension}\label{mod:IV-16-frame-at-every-end-as-the-stabiliser-of-its-dimension}
\chdesig{5-16}{4-16}{IV-16-frame-at-every-end-as-the-stabiliser-of-its-dimension}{IV-16-frame-at-every-end-as-the-stabiliser-of-its-dimension}
\begin{chwall}[The frame at every end as the stabiliser of its dimension]\label{box:IV-16-frame-at-every-end-as-the-stabiliser-of-its-dimension:1}
The frame at an end is the stabiliser of the pair ket of the link at that end, equal to the stabiliser of its dimension at the equal-end links and at the smaller end of an unequal link and parting at the larger end, the first end parting the qutrit end of the (2,3) pair: the stabiliser of the pair ket re-derived at the six cells from the face algebra has the dimensions four, four, seven, nine, nine and sixteen with the traceless parts three, three, six, eight, eight and fifteen and the centres one line, its projections to the two ends four and four, four and four, four and seven, nine and nine, nine and nine and sixteen and sixteen, and the stabiliser of its dimension at the three end dimensions is four, nine and sixteen with the traceless parts three, eight and fifteen; the Lorentz split at the qubit is re-met, the three rotation pairs inside the stabiliser of the (2,2) pair ket at three of three and equal to its traceless part, the three further generators carrying the pair ket; the vacuum stabiliser at the ququart is re-met, the (2,3) centre acting on the (4,4) register as the scalar zero and the (2,4) centre as minus two; the open result at the tau port is transcribed, the (2,2) port of the tau web paired to its leaf, the half rotation $-i,0;0,i$ with its square the full rotation minus one fixing the pair ket and the half-rotation generator inside the stabiliser of the (2,2) pair ket at one of one, an instance at the dimension two, the phrase at the qutrit not met at the port dimension two; the far faces of the six genres are defined at (2,2) under the transport reading and at (2,3) by the directions table and by no row at the four others, the far-end triple at the defined faces equal to the tensor-type row with the vector type at the Casimir two and the symmetric traceless type at the Casimir six, and against the stabiliser at the end the (2,2) far end gives four against four, equal, and the (2,3) far end, the qutrit, four against nine, parting. The recorded ends are three of seven, the tau port, the (2,2) far face and the (2,3) far face, equal at two of three with the frame object inside the stabiliser at every one; the first face, every frame an instance at its end dimension, reads no, and the second face reads yes and named, the qutrit end of the (2,3) pair at four against nine, the ququart ends of the (2,4) and (3,4) pairs parting beside at seven and nine against sixteen, the ends equal at four of seven; the open results of the frame row at their parting ends are equal at six of eleven, parting at three, the three at the hub at the larger end of its links, and at no end at two, every one recorded beside.
Re-posed at its own link: at the qutrit end of every one of the fifty-four links of the re-cut catalogue at the genre (2,3), each located at its own typed end at the dimension three, the frame, the stabiliser of the pair ket of that link projected to that end at the dimension four and carrying the far-end triple inside at three of three, parts from the stabiliser of the end dimension, the stabiliser of the pair ket at the equal-end cell of dimension three projected to the end at the dimension nine, four against nine at fifty-four of fifty-four links and equal at none, the end named the qutrit end of the first link of the first web in the catalogue order with both stabilisers printed, the same parting at every own link; the readings beside are printed, the support reading at the dimension two equal at four against four, the reduced reading parting at four against five and equal at five against five without the trace condition, the genre reading equal at four against four by construction; the stabiliser of the (2,3) pair ket is re-derived equal to the cells above at thirteen of thirteen and the first parting end above is re-met at four against nine.
\end{chwall}
\chstanding{open result}
\chdeps{Definitions 1.1 to 1.6, Theorems 9.1 to 9.7 and Open result 9.8, Theorem 54.1 and Theorem 124.1}{no theorem beyond it at this release}
The stabiliser of the pair ket is re-derived at the six cells from the face algebra with its projections to the ends; the stabiliser of each dimension is read; the Lorentz split, the vacuum stabiliser and the tau port open result are re-met at their cells; the far faces of the six genres are located at the tensor-type row and their triples re-derived; the stabiliser of the link projected to an end is compared with the stabiliser of that end's dimension at every end a recorded row defines a frame object at, and the first end parting is named.

\begin{chbridge}
The local symmetry at one end of a maximally entangled pair is the stabiliser of the pair, a proper subgroup of the unitary group of that end, and a frame at a point is the set of transformations fixing the state there \cite{Werner1989,Wigner1939}. The derivation supplies the stabilisers at every end and the first parting here, and the frame is named by correspondence.
\end{chbridge}
\begin{chproof}[finite and exact]
The premises are the face groups of the six cells (chapter 4-12, an instrument), the stabiliser of the pair ket (chapter 5-1, an open result), the Lorentz group at the qubit pair (chapter 3-3, a theorem), the electroweak breaking at the vacuum register (chapter 6-9, a theorem), the tensor type at every genre with its far faces (chapter 1-2, an open result) and the frame row at its eleven open results with the tau port (chapter 3-7, an open result). The chain back to the axiom runs through chapter 1-2 and the framework chapter. By exhaustive enumeration of the seven ends at which the comparison is posed, the tau port, the (2,2) far face, the (2,3) far face and the four far faces defined by no row, with the six cells at their stabilisers and the eleven open results of the frame row at their parting ends, the property tested at each end being whether the stabiliser of the pair ket projected to that end equals the stabiliser of the end's dimension. From the face algebra at a cell of end dimensions $f$ and $c$, by the action on the pair ket, we obtain the stabiliser
\begin{equation}
\dim \operatorname{Stab}(\Phi) = 4, 4, 7, 9, 9, 16 , \qquad \dim \operatorname{Stab}(d) = d^{2} = 4, 9, 16 ,
\end{equation}
the projections to the ends four and four at (2,2), four and four at (2,3), four and seven at (2,4), nine and nine at (3,3), nine and nine at (3,4) and sixteen and sixteen at (4,4), the central u(1) scalars zero, one, two, zero, one and zero inside the stabiliser at the equal-end cells alone, the finite face fixing the pair ket at sixteen, sixteen, sixty-four, sixty-four, sixty-four and two hundred fifty-six, equal to the face-groups table at six of six. From the (2,3) pair ket, by the projection to the qutrit end, the stabiliser carries the transported qubit triple with its centre at the dimension four against the stabiliser of the dimension three at nine,
\begin{equation}
4 \neq 9 \quad \text{at the qutrit end of the (2,3) pair} ,
\end{equation}
the first end parting, while at the (2,2) far end four against four is equal. The counts by kind: the Lorentz split re-met, the rotation pairs inside at three of three equal to the traceless part, the further generators carrying the ket with the first image $0,1,1,0$ equal to the Lorentz cell, the Gram and the ad matrices equal; the stabiliser at the qutrit re-met at the (2,3) cell at the dimension four with the centre $1;1;-1;-1;0$; the vacuum stabiliser re-met at the (2,3) centre as the scalar zero and the (2,4) centre as minus two at the genre-end reading, the stabiliser inside the span of dimension one at the direction $1;0$; the tau port open result at its located cell, the site of dimension two at degree one with its link at (2,2) and its far site of dimension two, the half rotation inside the stabiliser at one of one; the far faces defined at two of six with the triples equal to the tensor-type row at both and paired with the lifts inside the stabiliser at three of three; the ends equal at four of seven and the recorded ends at two of three; the open results of the frame row equal at six of eleven, the port open result, the dock open results at the qubit ends of (2,3) links, the dimension-as-valence open result at the qutrit end of its (3,4) links and the unit open result at the (4,4) word, parting at three of eleven at the hub at seven and nine against sixteen, and at no end at two. So the (2,3) far face where the tensor-type row found the type carries the qubit frame at the qutrit end and not the frame of the dimension three, which is the open result. The smallest instance on the page is the tau port: a qubit end at the (2,2) link, the half rotation squaring to minus one and fixing the pair ket, inside the stabiliser of dimension four.

The counter-tests read the port at the dimension three, which moved the first end to the port at four against nine. A second struck the transport identity, which moved the (2,2) far face inside the stabiliser from three of three to two of three. A third took the genre reading as the rule, which moved the deciding cell to yes with every end equal by construction, a fourth struck the gate, which left every cell struck, and a fifth put a product ket in place of the pair ket, which moved the (2,2) stabiliser dimension from four to two. Every control read the same at both engines, fourteen of fourteen, and the stabilisers were reproduced at a second engine of a different construction, byte-identical at every run.
\chcert{7}
\end{chproof}
\begin{figure}[h]\centering
\begin{adjustbox}{max width=\linewidth}\begin{tabular}{llrrl}\toprule
the end & its dimension & the stabiliser of the pair ket projected & the stabiliser of the dimension & equal \\ \midrule
the tau port, a qubit end of a (2,2) link & 2 & 4 & 4 & yes \\
the (2,2) far face under the transport reading & 2 & 4 & 4 & yes \\
the (2,3) far face by the directions table, the qutrit & 3 & 4 & 9 & no, the first parting \\
the ququart end of the (2,4) pair, beside & 4 & 7 & 16 & no \\
the ququart end of the (3,4) pair, beside & 4 & 9 & 16 & no \\
the ends of the (3,3) and (4,4) pairs, beside & 3 and 4 & 9 and 16 & 9 and 16 & yes \\ \bottomrule
\end{tabular}\end{adjustbox}
\caption{The stabiliser of the pair ket projected to an end against the stabiliser of that end's dimension at the seven ends, equal at the equal-end links and the smaller ends and parting at the larger end, the qutrit end of the (2,3) pair first.}
\end{figure}
